\documentclass[12pt]{article}
\usepackage{cite}
\usepackage{amsmath,amssymb,amsfonts,mathtools}
\usepackage{graphicx}
\usepackage{textcomp}
\usepackage{xcolor}
\usepackage{colortbl}
\usepackage{amsthm}
\usepackage[hidelinks]{hyperref}
\usepackage{caption}
\usepackage{array}
\usepackage{pgf,tikz,pgfplots}
\usepackage{mathrsfs}
\usepackage{lmodern}
\usepackage{fontspec}
\usepackage{listings}
\usepackage{longtable}
\usepackage{booktabs}
\usepackage{placeins}
\usepackage{enumitem}
\usetikzlibrary{arrows}
\pgfplotsset{compat=1.18}



\IfFontExistsTF{Menlo}
  {\newfontfamily{\leanlistingfont}{Menlo}[Scale=MatchLowercase]}
  {\IfFontExistsTF{DejaVu Sans Mono}
    {\newfontfamily{\leanlistingfont}{DejaVu Sans Mono}[Scale=MatchLowercase]}
    {\newfontfamily{\leanlistingfont}{Latin Modern Mono}[Scale=MatchLowercase]}}

\lstdefinelanguage{Lean4}{
  keywords={theorem, lemma, def, by, have, rw, simp, ring, norm_num, omega, let, show, from, fun, exact, refine, intro, rintro, revert, constructor, ext, import, open, section, end, where, match, with, if, then, else, do, return, variable, instance, class, structure, namespace, noncomputable, deriving},
  morecomment=[l]{--},
  morecomment=[n]{/-}{-/},
  morestring=[b]",
  sensitive=true,
}
\lstset{
  language=Lean4,
  basicstyle=\small\leanlistingfont,
  keywordstyle=\color{blue}\bfseries,
  commentstyle=\color{gray}\itshape,
  stringstyle=\color{red},
  showstringspaces=false,
  columns=fullflexible,
  keepspaces=true,
  breaklines=true,
  frame=single,
  backgroundcolor=\color{gray!10},
  captionpos=t,
  abovecaptionskip=0.8\baselineskip,
  belowcaptionskip=0.6\baselineskip,
  numbers=none,
  xleftmargin=1em,
  xrightmargin=1em,
  literate=
    {∀}{{$\forall$}}1
    {∃}{{$\exists$}}1
    {↔}{{$\leftrightarrow$}}1
    {→}{{$\to$}}1
    {←}{{$\leftarrow$}}1
    {≤}{{$\le$}}1
    {≥}{{$\ge$}}1
    {≠}{{$\neq$}}1
    {∈}{{$\in$}}1
    {∉}{{$\notin$}}1
    {⊆}{{$\subseteq$}}1
    {⊂}{{$\subset$}}1
    {∧}{{$\wedge$}}1
    {∨}{{$\vee$}}1
    {¬}{{$\neg$}}1
    {⊢}{{$\vdash$}}1
    {∑}{{$\sum$}}1
    {ℕ}{{$\mathbb{N}$}}1
    {ℤ}{{$\mathbb{Z}$}}1
    {ℚ}{{$\mathbb{Q}$}}1
    {ℝ}{{$\mathbb{R}$}}1
    {ℂ}{{$\mathbb{C}$}}1
    {𝔽}{{$\mathbb{F}$}}1
    {σ}{{$\sigma$}}1
    {μ}{{$\mu$}}1
    {η}{{$\eta$}}1
    {·}{{$\cdot$}}1
    {⟨}{{$\langle$}}1
    {⟩}{{$\rangle$}}1
    {ᵀ}{{$^{\mathsf{T}}$}}1
    {₀}{{$_0$}}1
    {₁}{{$_1$}}1
    {₂}{{$_2$}}1
    {₃}{{$_3$}}1
    {₄}{{$_4$}}1
    {₅}{{$_5$}}1
    {₆}{{$_6$}}1
    {₇}{{$_7$}}1
    {₈}{{$_8$}}1
    {₉}{{$_9$}}1
}
\captionsetup[lstlisting]{
  labelfont=bf,
  textfont=bf,
  singlelinecheck=false,
  justification=raggedright,
  skip=6pt
}

\usepackage{xpatch}
\makeatletter
\AtBeginDocument{\xpatchcmd{\@thm}{\thm@headpunct{.}}{\thm@headpunct{}}{}{}}
\makeatother

\setcounter{MaxMatrixCols}{20}
\setlength{\parskip}{3mm}
\setlength{\emergencystretch}{2em}
\topmargin=-30pt
\textheight=648pt
\oddsidemargin=0pt
\textwidth=468pt

\pagestyle{plain}
\newtheorem{defn}{Definition}[section]
\newtheorem{theorem}{Theorem}[section]
\newtheorem{lemma}{Lemma}[section]
\newtheorem{prop}[theorem]{Proposition}
\newtheorem{cor}[theorem]{Corollary}
\theoremstyle{definition}
\newtheorem{example}{Example}[section]
\theoremstyle{plain}
\newtheorem{question}[theorem]{Question}
\newtheorem{remark}{Remark}[section]

\newcommand{\bbF}{\mathbb{F}}
\newcommand{\F}{\mathbb{F}}
\newcommand{\Z}{\mathbb{Z}}
\newcommand{\Q}{\mathbb{Q}}
\newcommand{\N}{\mathbb{N}}
\newcommand{\GF}[1]{\ensuremath{\mathrm{GF}(#1)}}
\newcommand{\ZMod}[1]{\Z/#1\Z}
\newcommand{\ord}{\mathrm{ord}}
\newcommand{\Aut}{\mathrm{Aut}}
\newcommand{\Stab}{\mathrm{Stab}}
\newcommand{\Mon}{\mathrm{Mon}}
\newcommand{\St}{\mathrm{St}}
\newcommand{\rank}{\mathrm{rank}}
\newcommand{\Spec}{\mathrm{Spec}}
\newcommand{\Br}{\mathrm{Br}}
\newcommand{\wt}{\mathrm{wt}}
\newcommand{\Span}{\mathrm{Span}}
\newcommand{\et}{\mathrm{ét}}
\newenvironment{blueblock}{\begingroup\color{blue}}{\endgroup}

\begin{document}

\title{Formalizing building-up construction of self-dual codes through isotropic lines in Lean}

\author{
Jae-Hyun Baek \\
Department of Artificial Intelligence and \\
Institute for Mathematical and Data Sciences \\
Sogang University, Seoul, Korea
\and
Jon-Lark Kim\\
Department of Mathematics and\\
Institute for Mathematical and Data Sciences \\
Sogang University, Seoul, Korea
}

\date{August 27, 2026}
\maketitle

\begin{abstract}
We study building-up constructions of Euclidean self-dual codes through isotropic lines and formalize their algebraic structure in Lean~4. For \(q\equiv1\pmod4\), we prove that every self-dual code over \(\F_q\) has a repeated rank-\(r\) boxed generator after a permutation of two-coordinate blocks. The terminal \(r\times r\) matrix is nonsingular, and explicit Gram equations characterize self-duality. Simultaneous removal of pivot rows and coordinate blocks preserves this form; conversely, an explicit extension can be iterated while keeping the terminal matrix fixed. We determine precisely when such an extension is a Kim--Lee step over the same parent, and identify the complementary direct-sum case.
In the binary case, we give an independent formal proof of the Chinburg--Zhang rank-one boxed representation and compare its reduction with Kim's construction. A specified orthonormal-basis isometry identifies the arithmetic pairing with the Euclidean form; the distinguished Euclidean coordinate plane has Gram matrix \(I_2\), not the alternating hyperbolic Gram matrix. The arithmetic realization theorem remains a cited input.
Applications include optimal self-dual \([6,3,4]\) and \([8,4,4]\) codes over \(\GF{5}\), together with MDS self-dual codes over \(\GF{13}\).  In particular, a known pure double-circulant \([14,7,8]\) code is put into an explicit universal rank-one box; its two-coordinate reduction gives a self-dual \([12,6,6]\) parent with a unique projective correction capable of restoring distance~8.  A second exact reduction and reconstruction gives the sequence \([18,9,8]_{13}\leftrightarrow[20,10,10]_{13}\), attaining the published best-known distances at both lengths. Reproducible finite-field computations accompany these examples; they are separate from the kernel-checked algebraic proofs.
\end{abstract}




\section{Introduction}
\label{sec:intro}

Coding theory studies reliable communication, following Shannon's foundational work~\cite{Sha}. Among its classical examples, the extended binary Hamming code of length~8, the extended binary Golay code of length~24, and the extended ternary Golay code of length~12 are self-dual, that is, $C=C^\perp$.
Thus the classification or construction of self-dual codes has been one of the most active research problems~\cite{self-dual}.
Self-dual codes have wide connections with several mathematical areas including $t$-design theory~\cite{BacGab}, group theory~\cite{ConSlo}, number theory~\cite{BaDou},~\cite{ConSlo},~\cite{ShiChoShaSol}. Furthermore, they have applications in quantum error-correcting codes~\cite{CRSS}.

Kim~\cite{Kim2001} introduced ``building-up construction'' which proved that all binary self-dual codes can be constructed in this way. This construction starts from a binary self-dual code of length $2n$ and then adds a vector of length $2n$ of odd weight to construct a binary self-dual code of length $2n+2$. Kim--Lee~\cite{KimLee2009},~\cite{KimLee} generalized this method to construct self-dual codes over finite fields $\mathbb F_q$, where $q$ is even or $q \equiv 1 \pmod{4}$.
 In practice, this is a principal method for constructing examples over these finite fields. However, the formula for the new generator rows is often presented as a random vector whose inner product with itself is $-1$. One of the central points of the present paper is that a correct row vector is naturally organized by an isotropic line in a hyperbolic plane, and that the forward construction admits a precise row-wise and family-wise characterization in those terms.

Kreck and Puppe~\cite{KreckPuppe} realized binary self-dual codes using the cohomology of three-manifolds. Chinburg and Zhang~\cite{ChinburgZhang} obtained an arithmetic counterpart from the restriction map for \(\mu_2\), the sheaf of square roots of unity, on a ring of \(S\)-integers. Its image is self-annihilating for a bilinear pairing assembled from local Hilbert symbols. When the local pairings admit orthonormal bases, the image becomes a Euclidean binary self-dual code in those coordinates. Their arithmetic existence theorem, quoted as Theorem~\ref{thm-binary-hilbert-2}, realizes every binary self-dual code of length at least four up to permutation equivalence. We use this cited arithmetic input only through the restriction image, its pairing, and a chosen form isometry; these objects are specified before Theorem~\ref{thm:binary-cz-kim}. Our algebraic reduction and reconstruction do not require a formalization of arithmetic duality.

In this paper, we focus on $q \equiv 1 \pmod 4$. Over $\F_q$, self-dual codes are maximal totally isotropic subspaces for the standard Euclidean bilinear form on $\F_q^{2n}$. In this split case the elementary identity $c^2=-1$ is the basic algebraic input behind both the classical building-up construction and the structural description of the ambient quadratic space: it produces an isotropic line in $\F_q^2$, identifies the Euclidean plane with a hyperbolic plane, and makes the norm form $x^2+y^2$ split.


We establish that each lower generator row of the building-up admits the isotropic line decomposition $r_i = -y_i \cdot (1, c, 0, \ldots, 0) + (0, 0, g_i)$, where $(1, c)$ spans an isotropic line in the hyperbolic plane. This decomposition provides a precise geometric interpretation of the correction terms in the building-up construction. We also give an explicit representation by the quadratic form $x^2+y^2$: once $c^2=-1$ and $2\neq 0$, every $a\in\F_q$ can be written explicitly as $a=x^2+y^2$, which in turn yields extension vectors of prescribed norm.

The two arithmetic inputs must be distinguished. In the odd-characteristic $q$-ary part, $q \equiv 1 \pmod 4$ makes $-1$ a square in $\F_q^\times$, so the Euclidean plane is hyperbolic. In the binary construction of Chinburg--Zhang~\cite[Theorem~1.2 and Corollary~1.4]{ChinburgZhang}, the local cup-product pairing admits an orthonormal basis under the stated local hypotheses; for a non-complex place this is equivalent to $-1$ being a nonsquare in $K_v$. After the resulting coordinate isometry is fixed, the arithmetic image becomes a standard Euclidean self-dual binary code and its boxed matrix can be compared with Kim's construction.

Kim's binary building-up theorem and converse, the Kim--Lee extensions, and the Chinburg--Zhang arithmetic realization and binary boxed representation are established results. We distinguish them from the structural results proved here: the universal rank-\(r\) representation over split odd fields, its recursive restrictions and extensions, and its fixed-parent correspondence with Kim--Lee steps. The binary representation is given an independent algebraic proof suitable for formalization. The Lean development verifies these algebraic arguments and their interfaces; it does not formalize the cited arithmetic realization theorem.

Our main contributions are as follows.
\begin{enumerate}[label=(\arabic*)]



\item In the binary case, we separate the cohomological pairing from the standard Euclidean dot product by an explicit form isometry. In Euclidean coordinates the first coordinate plane is non-alternating with Gram matrix $I_2$; its diagonal vector $(1,1)$ is isotropic, and orthogonality forces the Kim correction formula. We prove intrinsically that every binary Euclidean self-dual code is permutation equivalent to the literal rank-one boxed form, by two-coordinate reduction and row normalization of the Kim extension (Theorem~\ref{thm:binary-universal-rank-one}). Deleting the first two coordinates and the top row recovers a shorter Euclidean self-dual code, and exact rebuilding recovers the boxed code (Theorem~\ref{thm:binary-cz-kim}).

\item In the $q$-ary case, we refine the classical building-up construction
into an exact self-duality criterion and prove
a universal rank-$r$ boxed normal form.  It is specified by terminal rows
\(\ell_i\), off-diagonal coefficients \(b_{ij}\), and a nonsingular
\(r\times r\) matrix \(D\); orthogonality forces the pivot diagonal and
the opposite block.  Simultaneous pivot-row and block-column
deletion preserves the form and self-duality
(Theorem~\ref{thm:universal-rank-boxed}); its rank-one specialization with
zero pivot diagonal recovers the explicit split boxed construction
(Corollary~\ref{thm:split-boxed-form}). Lemma~\ref{lem:rank-boxed-building-up}
constructs repeatable extensions with the same \(D\), and
Lemmas~\ref{lem:repeated-kim-lee}--\ref{lem:kim-lee-repeated}
identify their fixed-parent correspondence with Kim--Lee steps, including
the direct-sum exception.

\item As concrete applications, we exhibit the largest repeated
boxed presentations over \(\GF{2}\), \(\GF{5}\), and \(\GF{13}\), ending at
self-dual codes with parameters \([24,12,8]\), \([24,12,9]\), and
\([20,10,10]\), respectively.  For a known pure double-circulant
\([14,7,8]\) code over \(\GF{13}\), we also give an exact universal rank-one
generator, its \([12,6,6]\) two-coordinate parent, and the unique projective
correction that can restore distance~8 over this parent
(Section~\ref{sec:applications}).

\item We formalize the algebraic results in Lean~4, including the
self-duality criteria, binary two-coordinate reduction, the universal
normal forms, and the conditional coefficient relations. The formal
statements distinguish coordinate permutations from isometries between
different bilinear forms.
\end{enumerate}

\subsection*{Formalization in Lean}

The accompanying Lean~4 development formalizes the finite-dimensional
linear and bilinear algebra used in the paper. It includes the
self-dual/Lagrangian characterization, the systematic generator criterion,
the building-up constructions, both universal normal forms, the explicit
fixed-parent correspondence between repeated boxes and Kim--Lee steps, and the
conditional coefficient relations of
Theorem~\ref{thm:conditional-split-boxed-normalization}.
In the binary case, the proof of the universal normal form proceeds by
two-coordinate reduction and induction; it does not assume a classification
theorem. The quoted results of Kim, Kim--Lee, and Chinburg--Zhang remain
cited mathematical results and are not introduced as Lean axioms.
The finite-field searches and minimum-distance computations in the
applications are separate computations, not formalized proofs.

The Lean listings state the corresponding definitions and theorems; the
complete proofs and verification records are available at
\begin{center}
\href{https://github.com/LeGenAI/intersection-coding-theory-cohomology}{\texttt{LeGenAI/intersection-coding-theory-cohomology}}.
\end{center}
The development uses Lean~\texttt{v4.29.0-rc6} and Mathlib.
All 123 independently compared declarations are accepted by Lean's default
kernel. Their only foundational axioms are \texttt{propext},
\texttt{Quot.sound}, and \texttt{Classical.choice}; the proofs use neither
\texttt{sorry} nor \texttt{native\_decide}.
All seventeen suites were replayed together in a fresh Linux project
snapshot on August~27, 2026. All 123 declarations passed statement and
dependency comparison, the permitted-axiom audit, and default-kernel replay.
The recorded inputs were unchanged throughout the run. The numerical example
certificates and the external minimum-distance
computations also passed, but are not included in the 123-goal count.

The paper is organized as follows. Section~\ref{sec:background} fixes notation
and the split quadratic background. Section~\ref{sec:building-up} collects the
building-up tools and presents the universal rank-\(r\) boxed form, its binary
rank-one refinement, and the normalized rank-one forms over odd fields.
Section~\ref{sec:applications} gives the computational applications over
\(\GF{2}\), \(\GF{5}\), and \(\GF{13}\).
Section~\ref{sec:conclusion} closes with the comparison viewpoint and the
formalization perspective.


\section{Preliminaries}
\label{sec:background}


We refer to~\cite{HuffmanPless} for basic concepts and facts and to~\cite{self-dual} for self-dual codes.

Throughout the paper, \(\F_q\) denotes a finite field. We equip \(\F_q^n\) with the Euclidean bilinear form
\[
\langle u,v\rangle := \sum_{i=1}^n u_i v_i.
\]
Here ``Euclidean'' refers to the coding-theoretic dot product, not to
positive definiteness. We use \(\F_q\) for a general finite field and
\(\GF{5}\), \(\GF{13}\) for the concrete fields in the examples. For a
linear code, the minimum distance equals the minimum Hamming weight of
its nonzero words.
For a linear code $C \subseteq \F_q^n$, we write
\[
C^\perp := \{v \in \F_q^n : \langle v,c\rangle = 0 \text{ for all } c \in C\}.
\]

\begin{defn}[self-dual code]\label{def:sd}
{\rm
A linear code $C \subseteq \F_q^n$ is \emph{self-dual} if $C = C^{\perp}$~\cite{Pless1965}.
}
\end{defn}

\begin{defn}[Lagrangian subspace]\label{def:lagrangian}
{\rm
Let \(V\) be a finite-dimensional vector space over \(\F_q\) equipped with a nondegenerate symmetric bilinear form. A subspace \(L\subseteq V\) is called \emph{Lagrangian}~\cite{Mon,Ser} if \(L=L^\perp\).
}
\end{defn}

\noindent
For a nondegenerate form, \(\dim L^\perp=\dim V-\dim L\), so \(L=L^\perp\) if and only if \(L\) is totally isotropic and \(2\dim L=\dim V\). In particular, \(V\) is even-dimensional whenever it admits a Lagrangian subspace. Applied to the standard Euclidean form, this shows that a code \(C\subseteq\F_q^n\) is self-dual if and only if it is totally isotropic and \(2\dim C=n\). Thus a self-dual code is precisely a Lagrangian subspace of \((\F_q^n,\langle\cdot,\cdot\rangle)\), and its length is necessarily even.

\begin{lstlisting}[caption={Lean for Definitions~\ref{def:sd} and~\ref{def:lagrangian}.}]
def paperSelfDualCode (C : Submodule K (Fin n → K)) : Prop :=
  C = (dotBilin (K := K) (n := n)).orthogonal C

def paperLagrangianSubspace (L : Submodule K (Fin n → K)) : Prop :=
  L = (dotBilin (K := K) (n := n)).orthogonal L

theorem paper_self_dual_iff_lagrangian :
    paperSelfDualCode (K := K) C ↔ paperLagrangianSubspace (K := K) C

theorem paper_dotBilin_nondegenerate :
    (dotBilin (K := K) (n := n)).Nondegenerate

theorem paperLagrangianSubspace_iff_totallyIsotropic_and_finrank_half :
    paperLagrangianSubspace (K := K) L ↔
      IsTotallyIsotropic (dotBilin (K := K) (n := n)) L ∧
        2 * Module.finrank K L = Module.finrank K (Fin n → K)

theorem paperSelfDualCode_iff_totallyIsotropic_and_finrank_half :
    paperSelfDualCode (K := K) C ↔
      IsTotallyIsotropic (dotBilin (K := K) (n := n)) C ∧
        2 * Module.finrank K C = Module.finrank K (Fin n → K)

theorem paperLagrangianSubspace_even_length
    (hL : paperLagrangianSubspace (K := K) L) : Even n

theorem paperSelfDualCode_even_length
    (hC : paperSelfDualCode (K := K) C) : Even n
\end{lstlisting}

\noindent
For the odd-characteristic \(q\)-ary construction below, we further assume \(q\equiv 1\pmod 4\).

\begin{prop}[Systematic form criterion]\label{prop:systematic-form}
Let $C \subseteq \F_q^{2k}$ be generated by a matrix $G = [I_k \mid A]$ with $A \in M_k(\F_q)$. Then $C$ is self-dual if and only if
\[
AA^T = -I_k.
\]
\end{prop}

\begin{proof}
The rows of $G$ span a $k$-dimensional code. Hence self-duality is equivalent to pairwise orthogonality of the rows. The Gram matrix $GG^T$ of the rows is
\[
GG^T = I_k + AA^T = O,
\]
so the orthogonality condition is exactly $AA^T = -I_k$.
\end{proof}

\begin{lstlisting}[caption={Lean for Proposition~\ref{prop:systematic-form}.}]
def systematicRows (A : Matrix (Fin k) (Fin k) K) :
    Fin k → Fin (k + k) → K :=
  fun i => Fin.append (Pi.single i 1) (A i)

theorem systematicRows_linearIndependent (A : Matrix (Fin k) (Fin k) K) :
    LinearIndependent K (systematicRows A)

theorem paper_systematic_form_criterion_exact
    (A : Matrix (Fin k) (Fin k) K) :
    paperSelfDualCode (rowSpace (systematicRows A)) ↔
      A * A.transpose = -(1 : Matrix (Fin k) (Fin k) K)
\end{lstlisting}

\begin{defn}[Permutation Equivalence]\label{def:perm-equiv}
{\rm
Two codes $C_1,C_2\subseteq\F_q^n$ are \emph{permutation equivalent} if there is a permutation $\sigma\in S_n$ such that $C_2=P_\sigma C_1$, where $(P_\sigma u)_j=u_{\sigma(j)}$.
}
\end{defn}

Coordinate permutations are Euclidean isometries:
\[
\langle P_\sigma u,P_\sigma v\rangle=\langle u,v\rangle.
\]
Consequently, permutation equivalence preserves Euclidean self-duality. A general monomial transformation additionally permits independent nonzero coordinate scalings; it need not preserve the Euclidean form or self-duality and is not included in our use of ``equivalence'' below. Over $\F_2$, permutation and monomial equivalence coincide because $\F_2^\times=\{1\}$.

\begin{lstlisting}[caption={Lean for Definition~\ref{def:perm-equiv}.}]
def CodeEquiv (G1 : Fin m₁ → Fin n → K)
    (G2 : Fin m₂ → Fin n → K) : Prop :=
  ∃ σ : Equiv.Perm (Fin n),
    rowSpace (permuteFamily σ G1) = rowSpace G2

theorem dot_coordinatePermuteLinearEquiv :
    dot (permuteVec σ u) (permuteVec σ v) = dot u v

theorem codeEquiv_preserves_paperSelfDualCode
    (hG1G2 : CodeEquiv G1 G2)
    (hG1 : paperSelfDualCode (K := K) (rowSpace G1)) :
    paperSelfDualCode (K := K) (rowSpace G2)
\end{lstlisting}





From this point onward, when discussing the split \(q\)-ary construction we assume that \(q\) is odd and satisfies \(q \equiv 1 \pmod 4\). The condition \(q \equiv 1 \pmod 4\) is equivalent to $-1$ being a square in $\F_q^\times$. We fix once and for all an element
\[
c \in \F_q^\times \qquad \text{with} \qquad c^2 = -1.
\]
This choice is used throughout the building-up formulas.


\begin{defn}[Hyperbolic Plane]\label{def:hyp}
{\rm
A \emph{hyperbolic plane}~\cite{Ser} over $\F_q$ is a $2$-dimensional symmetric bilinear space with basis $\{e_1,e_2\}$ satisfying
\[
\langle e_1,e_1\rangle = \langle e_2,e_2\rangle = 0,
\qquad
\langle e_1,e_2\rangle = 1.
\]
Such a basis \(\{e_1,e_2\}\) will be called a \emph{hyperbolic pair}. Any $1$-dimensional totally isotropic subspace is called an \emph{isotropic line}.
}
\end{defn}

In even dimension and odd characteristic, a nondegenerate symmetric bilinear
space is called \emph{split} if it is an orthogonal direct sum of hyperbolic
planes. Thus the splitting condition concerns the bilinear space, not the
underlying finite field.

\begin{prop}\label{prop:split-consequences}
Let \(q\) be odd and let \(c\in\F_q\) satisfy \(c^2=-1\). Then:
\begin{enumerate}[label=(\roman*)]
\item The element $c$ has multiplicative order $4$.
\item The Euclidean plane $(\F_q^2,\langle \cdot,\cdot\rangle)$ is hyperbolic.
\item For every $k \ge 1$, the matrix equation $AA^T = -I_k$ has the explicit solution $A=cI_k$.
\item The quadratic form $x^2+y^2$ splits as
\[
x^2+y^2 = (x-cy)(x+cy),
\]
hence it represents every element of\/ $\F_q$.
\end{enumerate}
\end{prop}

\begin{proof}
Since $q$ is odd, $-1\ne 1$.  From $c^2=-1$ we obtain
\[
c^4=(c^2)^2=(-1)^2=1,
\qquad c^2=-1\ne 1.
\]
Thus the order of $c$ divides $4$, but does not divide $2$, and
therefore $\operatorname{ord}(c)=4$.  Put
\[
e_1=(1,c),\qquad e_2=(2^{-1},-2^{-1}c).
\]
Then
\[
\langle e_1,e_1\rangle=1+c^2=0,
\qquad
\langle e_2,e_2\rangle=2^{-2}(1+c^2)=0,
\]
whereas
\[
\langle e_1,e_2\rangle
=2^{-1}(1-c^2)=2^{-1}(1-(-1))=1.
\]
If $\alpha e_1+\beta e_2=0$, taking inner products first with
$e_2$ and then with $e_1$ gives
\[
\alpha=\langle \alpha e_1+\beta e_2,e_2\rangle=0,
\qquad
\beta=\langle \alpha e_1+\beta e_2,e_1\rangle=0.
\]
Hence $e_1,e_2$ are linearly independent and consequently form a
hyperbolic basis of the $2$-dimensional Euclidean plane.

For $A=cI_k$, symmetry of $I_k$ and $c^2=-1$ give
\[
AA^T=(cI_k)(cI_k)^T=c^2I_k=-I_k.
\]
Finally,
\[
(x-cy)(x+cy)=x^2-c^2y^2=x^2+y^2.
\]
For an arbitrary $a\in\F_q$, take
\[
x=\frac{1+a}{2},
\qquad
y=\frac{1-a}{2}\,c.
\]
Since $c^2=-1$, we have
\[
cy=-\frac{1-a}{2},
\qquad x-cy=1,
\qquad x+cy=a,
\]
and hence $x^2+y^2=(x-cy)(x+cy)=a$.
\end{proof}


\begin{prop}\label{prop:eucl-hyp}
Let \(q\) be odd and let \(c\in\F_q\) satisfy \(c^2=-1\). Then the vectors
\[
e_1=(1,c), \qquad e_2=(2^{-1},-2^{-1}c)
\]
form a hyperbolic basis of\/ $\F_q^2$. Equivalently, the line spanned by $(1,c)$ is isotropic and the Euclidean plane is split.
\end{prop}

\begin{proof}
The proof of Proposition~\ref{prop:split-consequences}(ii) computes
\(\langle e_i,e_i\rangle=0\), \(\langle e_1,e_2\rangle=1\), and
establishes linear independence of these two vectors. These are exactly
the required basis and Gram conditions.
\end{proof}

\begin{lstlisting}[caption={Lean for Proposition~\ref{prop:split-consequences} and Proposition~\ref{prop:eucl-hyp}.}]
theorem paper_split_consequences_exact
    [Fact ((2 : K) ≠ 0)] (c : K) (hc : c ^ 2 = (-1 : K)) :
    orderOf c = 4 ∧
      (paperHyperbolicPair (splitE1 c) (splitE2 c) ∧
        LinearIndependent K ![splitE1 c, splitE2 c]) ∧
      (∀ k : Nat,
        (c • (1 : Matrix (Fin k) (Fin k) K)) *
            (c • (1 : Matrix (Fin k) (Fin k) K)).transpose =
          -(1 : Matrix (Fin k) (Fin k) K)) ∧
      (∀ a : K, ∃ x y : K,
        x - c * y = 1 ∧ x + c * y = a ∧ x ^ 2 + y ^ 2 = a)

theorem norm_form_witness_factors
    [Fact ((2 : K) ≠ 0)] (a c : K) (hc : c ^ 2 = (-1 : K)) :
    let x := (1 + a) / 2
    let y := ((1 - a) / 2) * c
    x - c * y = 1 ∧ x + c * y = a ∧ x ^ 2 + y ^ 2 = a

theorem paper_euclidean_plane_hyperbolic_basis_exact
    [Fact ((2 : K) ≠ 0)] (c : K) (hc : c ^ 2 = (-1 : K)) :
    paperHyperbolicPair (splitE1 c) (splitE2 c) ∧
      LinearIndependent K ![splitE1 c, splitE2 c]
\end{lstlisting}

A change to hyperbolic coordinates is distinct from permutation equivalence.
Suppose \(2\ne0\) and \(c^2=-1\), and write \(v=(v_1,v_2,v')\).
The invertible linear map
\[
A_c(v)=\left(\frac{v_1-cv_2}{2},\,\frac{v_1+cv_2}{2},\,v'\right)
\]
satisfies \(B_{\rm sp}(A_cu,A_cv)=\langle u,v\rangle\), where
\[
B_{\rm sp}(u,v)=2(u_1v_2+u_2v_1)+u'\cdot v'.
\]
Consequently, \(C=C^{\perp_{\rm E}}\) if and only if
\(A_cC=(A_cC)^{\perp_{B_{\rm sp}}}\).
The map \(A_c\) does not preserve the Euclidean form: \(A_c(1,c,0)=(1,0,0)\),
whose self-pairing is \(1\), whereas \((1,c,0)\) has self-pairing \(0\).
Thus \(A_c\) is not a Euclidean isometry and cannot replace a coordinate
permutation in our normal-form theorems.

\section{Building-Up and Boxed Normal Forms}
\label{sec:building-up}

We first collect the building-up constructions and the self-duality criterion
used below.  We then organize the boxed forms by their normalization:
the universal rank-\(r\) pattern in
Section~\ref{sec:universal-boxed}, the binary rank-one refinement in
Section~\ref{sec:binary-refinement}, and the normalized rank-one construction
over odd fields in Section~\ref{sec:odd-rank-one}.  Their computational
realizations are collected separately in Section~\ref{sec:applications}.

\subsection{Building-up constructions}
\label{sec:building-up-tools}

In 2001, Kim~\cite{Kim2001} proposed a building-up construction of binary self-dual codes in order to construct many self-dual codes of larger lengths from a self-dual code of a given length. An explicit form of the building-up construction is given below.

\begin{theorem}(\cite[Theorem 1]{Kim2001}) \label{thm-binary-build-1}
Let $G_0$ be a generator matrix of a binary self-dual code $C_0$ of length $2n$, whose rows are $g_1, \dots, g_n$. Let $\bf x$ be a binary vector in $\mathbb F_2^{2n}$ satisfying ${\bf x} \cdot {\bf x} = 1$. Suppose that $y_i ={\bf x} \cdot g_i$ for $1 \le i \le n$. Then the following matrix
\[
\begingroup
\setlength{\arraycolsep}{6pt}
\renewcommand{\arraystretch}{1.2}
G=
\left[
\begin{array}{cc!{\vrule width 1.2pt}c}
1 & 0 & \mathbf{x}\\
\noalign{\hrule height 1.2pt}
y_1 & y_1 & g_1\\
\vdots & \vdots & \vdots\\
y_n & y_n & g_n
\end{array}
\right]
\endgroup
\]
generates a binary self-dual code $C$ of length $2n+2$.
\end{theorem}

\begin{proof}
Write
\[
r_0=(1,0,\mathbf{x}),
\qquad
r_i=(y_i,y_i,g_i)\quad(1\le i\le n).
\]
Because $C_0$ is self-dual, its generator rows $g_1,\ldots,g_n$
are linearly independent and satisfy $g_i\mathbin{\cdot}g_j=0$ for
all $i,j$.  In characteristic two,
\[
\begin{aligned}
r_0\mathbin{\cdot}r_0
  &=1+\mathbf{x}\mathbin{\cdot}\mathbf{x}=1+1=0,\\
r_0\mathbin{\cdot}r_i
  &=y_i+\mathbf{x}\mathbin{\cdot}g_i=y_i+y_i=0,\\
r_i\mathbin{\cdot}r_j
  &=2y_i y_j+g_i\mathbin{\cdot}g_j=0.
\end{aligned}
\]
Thus the rows of $G$ are pairwise orthogonal.

It remains to verify their independence.  Suppose
\[
\lambda_0r_0+\sum_{i=1}^{n}\lambda_i r_i=0.
\]
Subtracting the second coordinate of this relation from the first gives
$\lambda_0=0$.  The tail coordinates then give
\[
\sum_{i=1}^{n}\lambda_i g_i=0,
\]
so $\lambda_i=0$ for every $i$ by the independence of the rows of
$G_0$.  Hence $G$ has $n+1$ independent, mutually orthogonal rows in
$\mathbb F_2^{2n+2}$.  Its row space is therefore totally isotropic of
dimension $n+1=\frac12(2n+2)$, and consequently equals its orthogonal
complement.
\end{proof}

\noindent\begin{minipage}{\linewidth}
\begin{lstlisting}[caption={Exact Lean theorem for Theorem~\ref{thm-binary-build-1}.}]
theorem paper_binary_kim_building_up_exact [CharP K 2]
    {m : Nat} {x : Fin (2 * m) → K}
    {G : Fin m → Fin (2 * m) → K}
    (hx : dot x x = (1 : K))
    (hparent : paperSelfDualCode (rowSpace G)) :
    paperSelfDualCode (rowSpace (buildRowsBin x G))
\end{lstlisting}
\end{minipage}

Then the converse of Theorem~\ref{thm-binary-build-1} was proved in~\cite{Kim2001} as follows.

\begin{theorem} (\cite[Theorem 2]{Kim2001})  \label{thm-binary-build-2}
 Any binary self-dual code $C$ of length $2n$ with minimum distance $d>2$ is obtained from some self-dual code $C_0$ of length
$2n-2$ (up to permutation equivalence) by the construction in Theorem~\ref{thm-binary-build-1}.
\end{theorem}

For the odd-characteristic constructions in the rest of this subsection,
write \(K=\F_q\) with \(q\equiv1\pmod4\).

Motivated by the building-up construction of binary self-dual codes, Kim and Lee~\cite{KimLee} generalized it to finite fields with $q$ elements where $q \equiv 1 \pmod{4}$ as follows.

\begin{theorem}(\cite[Proposition 2.1]{KimLee}) \label{thm-q-ary-build-1}
Let $\mathbb F_q$ be a finite field with $q$ elements. Assume that $q \equiv 1 \pmod{4}$. Let $c$ be in $\mathbb F_q$ such that $c^2 = -1$. Let $G_0$ be a generator matrix of a Euclidean self-dual code $C_0$ over $\mathbb F_q$ of length $2n$, whose rows are $g_1, \dots, g_n$. Let $\bf x$ be a vector in $\mathbb F_q^{2n}$ satisfying ${\bf x} \cdot {\bf x} = -1$. Suppose that $y_i ={\bf x} \cdot g_i$ for $1 \le i \le n$. Then the following matrix
\[
\begingroup
\setlength{\arraycolsep}{6pt}
\renewcommand{\arraystretch}{1.2}
G=
\left[
\begin{array}{cc!{\vrule width 1.2pt}c}
1 & 0 & \mathbf{x}\\
\noalign{\hrule height 1.2pt}
-y_1 & c y_1 & g_1\\
\vdots & \vdots & \vdots\\
-y_n & c y_n & g_n
\end{array}
\right]
\endgroup
\]
generates a self-dual code $C$ over $\mathbb F_q$ of length $2n+2$.
\end{theorem}

\begin{proof}
Let
\[
r_0=(1,0,\mathbf{x}),
\qquad
r_i=(-y_i,c y_i,g_i)\quad(1\le i\le n).
\]
Since $C_0=C_0^\perp$, the rows of $G_0$ are linearly independent and
$g_i\cdot g_j=0$ for all $i,j$.  Using $y_i=\mathbf{x}\cdot g_i$ and
$c^2=-1$, we obtain
\[
\begin{aligned}
r_0\cdot r_0
  &=1+\mathbf{x}\cdot\mathbf{x}=1-1=0,\\
r_0\cdot r_i
  &=-y_i+\mathbf{x}\cdot g_i=0,\\
r_i\cdot r_j
  &=(-y_i)(-y_j)+(c y_i)(c y_j)+g_i\cdot g_j\\
  &=(1+c^2)y_i y_j+g_i\cdot g_j=0.
\end{aligned}
\]
Thus the row space of $G$ is self-orthogonal.

For linear independence, suppose
\[
\lambda_0r_0+\sum_{i=1}^{n}\lambda_i r_i=0.
\]
The second coordinate gives
\[
c\sum_{i=1}^{n}\lambda_i y_i=0.
\]
Now $c\ne0$, since $c^2=-1$, and hence
$\sum_i\lambda_i y_i=0$.  The first coordinate then gives
$\lambda_0-\sum_i\lambda_i y_i=\lambda_0=0$.  Taking tail coordinates
yields
\[
\sum_{i=1}^{n}\lambda_i g_i=0,
\]
so every $\lambda_i$ vanishes.  Consequently $G$ has $n+1$ independent
pairwise orthogonal rows in $\mathbb F_q^{2n+2}$.  Its row space has
dimension
\[
n+1=\frac{2n+2}{2},
\]
and therefore equals its Euclidean orthogonal complement.
\end{proof}

\noindent\begin{minipage}{\linewidth}
\begin{lstlisting}[caption={Exact Lean theorem for Theorem~\ref{thm-q-ary-build-1}.}]
theorem paper_qary_kim_lee_building_up_exact
    {m : Nat} {x : Fin (2 * m) -> K} {c : K}
    {G : Fin m -> Fin (2 * m) -> K}
    (hx : dot x x = (-1 : K))
    (hc : c ^ 2 = (-1 : K))
    (hparent : paperSelfDualCode (rowSpace G)) :
    paperSelfDualCode (rowSpace (buildRows x (-c) G))
\end{lstlisting}
\end{minipage}

\begin{remark}
\label{rem:kim-lee-sign-convention}
\rm
The sign convention of Kim--Lee in the displayed matrix is $(-y_i,c y_i,g_i)$.
From this point on we use the equivalent normalization with rows
$(-y_i,-c y_i,g_i)$. Since
$(-c)^2=c^2=-1$, this changes only the sign convention in the second
correction coordinate.
\end{remark}

 Then the converse of the above theorem was also stated and proved in \cite{KimLee}.

\begin{theorem}(\cite[Proposition 2.2]{KimLee}) \label{thm-q-ary-build-2}
Let $\mathbb F_q$ be a finite field with $q$ elements. Assume that $q \equiv 1 \pmod{4}$. Any Euclidean self-dual code $C$ over $\mathbb F_q$ of length $2n$ with minimum distance $d>2$ is obtained from some Euclidean self-dual code $C_0$ over $\mathbb F_q$ of length $2n-2$ (up to permutation equivalence) by the construction method in Theorem~\ref{thm-q-ary-build-1}.
\end{theorem}

Kim and Lee define two
linear codes to be permutation equivalent when a permutation of coordinates
sends one code to the other, and Proposition~2.2 uses the phrase ``up to
permutation equivalence''~\cite[pp.~81, 83--84]{KimLee}.  The row equivalence
appearing inside their proof only changes the chosen generator basis and does
not enlarge the code equivalence relation to monomial equivalence.

Before imposing a recursive boxed shape, we determine precisely when the
building-up matrix generates a self-dual code, without a minimum-distance
assumption.

The ambient splitting used below is independent of the code and of its
generator matrix.  In the literal standard-coordinate decomposition
\[
W=E_2\perp W_0,
\qquad E_2=K^2,\qquad W_0=K^{2n-2},
\]
the plane $E_2$ has Gram matrix $I_2$.  Since the characteristic is odd, put
\[
e=(1,c),
\qquad
f=\left(\frac12,-\frac c2\right).
\]
Then
\[
e\cdot e=1+c^2=0,
\qquad
f\cdot f=\frac{1+c^2}{4}=0,
\qquad
e\cdot f=\frac{1-c^2}{2}=1.
\]
Thus $(e,f)$ is a hyperbolic basis of $E_2$, and the basis-change map gives
the form isometry $E_2\simeq H$.  Consequently $W\simeq H\perp W_0$,
whereas the matrix below remains written in the original standard Euclidean
coordinates.  This distinction is precisely Proposition~\ref{prop:eucl-hyp};
it is not inferred from the displayed generator matrix.

The next theorem combines the forward and reverse calculations.
No prescribed block structure is imposed on \(G\).
The theorem determines exactly when adjoining the displayed row and two
columns to \(G\) produces a generator matrix of a self-dual code.

\begin{theorem}[Self-duality criterion for the building-up matrix]\label{thm:qary-building-up}
Let $q$ be a prime power with $q \equiv 1 \pmod{4}$, and let $K=\mathbb{F}_q$.
Fix $c \in K^\times$ such that $c^2=-1$.
Let \(x\in K^{2n-2}\), let \(y_1,\dots,y_{n-1}\in K\), and let
\(G\) be an arbitrary \((n-1)\times(2n-2)\) matrix with rows
\(g_1,\dots,g_{n-1}\).  Form the bordered matrix
\[
\begingroup
\setlength{\arraycolsep}{6pt}
\renewcommand{\arraystretch}{1.2}
G'=
\left[
\begin{array}{cc!{\vrule width 1.2pt}c}
1 & 0 & x\\
\noalign{\hrule height 1.2pt}
-y_1 & -c y_1 & g_1\\
\vdots & \vdots & \vdots\\
-y_{n-1} & -c y_{n-1} & g_{n-1}
\end{array}
\right]
\endgroup,
\]
and put \(L=\operatorname{rowspan}(G')\) and
\(L_0=\operatorname{rowspan}(G)\subseteq K^{2n-2}\).  Then
\[
L=L^\perp
\quad\Longleftrightarrow\quad
\left\{
\begin{aligned}
&x\cdot x=-1,\\
&y_i=x\cdot g_i &&(1\le i\le n-1),\\
&L_0=L_0^\perp.
\end{aligned}
\right.
\]
Whenever these equivalent conditions hold, \(G'\) is literally the normalized
Kim--Lee building-up matrix obtained from \(G\) and \(x\), not merely a
generator of an equivalent code. Deleting its first row and first two
columns recovers \(G\) exactly.
\end{theorem}

\begin{proof}
Suppose first that \(L=L^\perp\).  Since the Euclidean form on \(K^{2n}\) is
nondegenerate, \(\dim L=n\).  The \(n\) displayed rows therefore form a basis
of \(L\).  Write
\[
r_0 = (1,0,x), \qquad r_i = (-y_i,-c y_i,g_i).
\]
Self-orthogonality gives
\[
0=\langle r_0,r_0\rangle=1+\langle x,x\rangle,
\]
so $\langle x,x\rangle=-1$. Likewise,
\[
0 = \langle r_0,r_i\rangle=-y_i+\langle x,g_i\rangle,
\]
and hence \(y_i=\langle x,g_i\rangle\) for every \(i\).  For
\(1\le i,j\le n-1\),
\[
\langle r_i,r_j\rangle
= y_i y_j + c^2 y_i y_j + \langle g_i,g_j\rangle
= (1 + c^2)y_i y_j + \langle g_i,g_j\rangle.
\]
Since $c^2 = -1$, the first term vanishes, and thus
\[
\langle r_i,r_j\rangle = \langle g_i,g_j\rangle.
\]
Because \(L\) is self-orthogonal, this proves
\(\langle g_i,g_j\rangle=0\) for all \(i,j\), so \(L_0\) is
self-orthogonal.

To determine its dimension, suppose
\[
\sum_{i=1}^{n-1}a_i g_i=0.
\]
The coefficient formula gives
\[
\sum_{i=1}^{n-1}a_i y_i
=\left\langle x,\sum_{i=1}^{n-1}a_i g_i\right\rangle=0.
\]
Consequently all three blocks of $\sum_i a_i r_i$ vanish.  The displayed
rows are linearly independent, so all $a_i$ vanish.  Thus the tail rows are
linearly independent and
\[
\dim L_0=n-1=\frac12\dim W_0.
\]
The self-orthogonality of \(L_0\), together with this half-dimension equality
and nondegeneracy of the Euclidean form on \(K^{2n-2}\), gives
\(L_0=L_0^\perp\).

Conversely, assume \(x\cdot x=-1\), \(y_i=x\cdot g_i\) for every \(i\), and
\(L_0=L_0^\perp\). Substitution of the coefficient identities gives
the rows \((1,0,x)\) and
\((-x\cdot g_i,-c(x\cdot g_i),g_i)\).
Theorem~\ref{thm-q-ary-build-1}, applied with the square root \(-c\),
therefore shows that \(L\) is self-dual.  Finally, tail projection
of every row other than the first deletes the first two coordinates and returns \(g_i\),
so deletion recovers \(G\) exactly.
\end{proof}



\noindent
In the following listing, \texttt{qaryFreeCoreBoxedFamily x c Y G} denotes
the displayed bordered matrix, with arbitrary \(G\). The equivalence states
the norm and coefficient conditions and the self-duality of the row space
of \(G\); the final equalities express the building-up and deletion formulas.

\noindent\begin{minipage}{\linewidth}
\begin{lstlisting}[caption={Lean for Theorem~\ref{thm:qary-building-up}.}]
theorem paper_qary_free_core_boxed_equivalence
    (hc : c ^ 2 = (-1 : K)) :
    paperSelfDualCode
      (rowSpace (qaryFreeCoreBoxedFamily x c Y G)) ↔
    dot x x = (-1 : K) ∧
    (∀ i, Y i = dot x (G i)) ∧
    paperSelfDualCode (rowSpace G) ∧
    qaryFreeCoreBoxedFamily x c Y G = buildRows x c G ∧
    deleteHyperbolicPairSplit
      (qaryFreeCoreBoxedFamily x c Y G) = G
\end{lstlisting}
\end{minipage}

\begin{remark}
\label{rem:ambient-form-and-generator}
{\rm
Proposition~\ref{prop:eucl-hyp} identifies the Euclidean plane with a
hyperbolic plane independently of the generator matrix.
Theorem~\ref{thm:qary-building-up} characterizes self-duality for arbitrary
\(G\); it does not impose a boxed normal form on \(G\). We turn to such
normal forms next.
}
\end{remark}

\subsection{Universal rank-\texorpdfstring{$r$}{r} boxed normal form}
\label{sec:universal-boxed}

Theorem~\ref{thm:qary-building-up} characterizes self-duality of the displayed
bordered matrix without prescribing a block structure for \(G\).
Under its equivalent conditions, the matrix is exactly the normalized
Kim--Lee building-up matrix.  In the minimum-distance range
\(d>2\), Kim--Lee's Proposition~2.2 further shows that every \(q\)-ary
self-dual code is permutation equivalent to such a presentation.  The next
theorem specifies a recursive block structure for the generator matrix,
designed to play for \(q\equiv 1\pmod 4\) the same recursive role that boxed
matrices play in the binary theorem of Chinburg--Zhang.  In odd characteristic
the terminal part need not consist of a single block column, and the pivot
diagonal need not be fixed to \(01\).  We first retain both freedoms to obtain
a universal, recursively closed form, and then examine its binary refinement and normalized rank-one specialization
in the following subsections.

For a terminal row
\(\ell=((u_t,v_t))_{t=1}^r\in(K^2)^r\), put
\[
\partial_c\ell=(v_t-cu_t)_{t=1}^r{}^{\mathsf T}\in K^r,
\qquad
\ell\mathbin{\boldsymbol\cdot}\ell'
=\sum_{t=1}^r (u_tu'_t+v_tv'_t).
\]
Thus \(\partial_c\ell\) records the transverse coordinates of \(\ell\),
while \(\boldsymbol\cdot\) is simply the Euclidean product on its \(2r\)
entries.  These are the only two operations on terminal rows needed below.

\begin{theorem}[Universal rank-\texorpdfstring{$r$}{r} boxed normal form over
\texorpdfstring{$\F_q$}{Fq}]
\label{thm:universal-rank-boxed}
Let \(q\equiv1\pmod4\), let \(K=\F_q\), and fix \(c\in K^\times\) with
\(c^2=-1\), and identify
\(K^{2n}=(K^2)^n\).  For every Euclidean self-dual code
\(C\subseteq K^{2n}\), define
\[
U_c\coloneqq\bigoplus_{j=1}^n K(1,c),
\qquad
r\coloneqq\dim(C\cap U_c),
\qquad k\coloneqq n-r.
\]
Then, after a permutation of the \(n\) two-coordinate blocks and elementary
row operations, \(C\) has a generator matrix of the following repeated form.
Its first \(k\) rows are the pivot rows, and its last \(r\) rows form a basis
of \(C\cap U_c\).  Choose terminal rows
\(\ell_1,\ldots,\ell_k\in(K^2)^r\), coefficients \(b_{ij}\in K\)
for \(i\ne j\), and put
\[
a_i=\frac c2\bigl(1+\ell_i\mathbin{\boldsymbol\cdot}\ell_i\bigr).
\]
Then \(C\) is generated by \(G(c;b,\ell,D)\), where
\begin{equation}
\label{eq:universal-rank-box}
\begingroup
\setlength{\arraycolsep}{2pt}
\renewcommand{\arraystretch}{1.35}
G(c;b,\ell,D)=
\left[
\begin{array}{cccc|c}
(a_1,ca_1+1) & b_{12}(1,c) & \cdots & b_{1k}(1,c) & \ell_1\\
b_{21}(1,c) & (a_2,ca_2+1) & \cdots & b_{2k}(1,c) & \ell_2\\
\vdots & \vdots & \ddots & \vdots & \vdots\\
b_{k1}(1,c) & b_{k2}(1,c) & \cdots & (a_k,ca_k+1) & \ell_k\\
\hline
-\bigl(D\partial_c\ell_1\bigr)_{s}(1,c) &
-\bigl(D\partial_c\ell_2\bigr)_{s}(1,c) & \cdots &
-\bigl(D\partial_c\ell_k\bigr)_{s}(1,c)
  & \bigl(D_{st}(1,c)\bigr)
\end{array}
\right],
\endgroup
\end{equation}
where the bottom line is indexed by \(1\le s\le r\),
and \(D\in M_r(K)\) satisfies
\begin{equation}
\label{eq:universal-rank-core}
\det D\ne0.
\end{equation}
The only remaining relations are
\begin{equation}
 c(b_{ij}+b_{ji})
 +\ell_i\mathbin{\boldsymbol\cdot}\ell_j=0
 \qquad(i\ne j).
 \label{eq:universal-rank-pp}
\end{equation}
Conversely, for every \(k,r\geq0\), every matrix
\(G(c;b,\ell,D)\) satisfying
\eqref{eq:universal-rank-core} and \eqref{eq:universal-rank-pp} has
\(k+r\) linearly independent, pairwise orthogonal rows and therefore
generates a Euclidean self-dual code of length \(2(k+r)\).

Moreover, the form is closed under simultaneous deletion of any pivot rows
and their corresponding block columns.  More precisely, for any ordered
subset \(I\subseteq\{1,\ldots,k\}\) of size \(l\), the retained matrix is
\begin{equation}
\label{eq:rank-boxed-restriction}
G_I=G(c;b|_{I\times I},(\ell_i)_{i\in I},D),
\end{equation}
and it generates a self-dual code of length \(2(l+r)\), with the same terminal matrix
\(D\).  At \(I=\varnothing\), its row space is exactly
\(U_c^{(r)}\coloneqq\{(z_1(1,c),\ldots,z_r(1,c)):z\in K^r\}\).

\end{theorem}

\begin{proof}
For \(v=(v_1,\dots,v_n)\in(K^2)^n\), write each block uniquely as
\[
v_j=\bigl(\alpha_j(v),c\alpha_j(v)+\beta_j(v)\bigr),
\qquad
\alpha_j(v)=v_{j,1},
\qquad
\beta_j(v)=v_{j,2}-cv_{j,1}.
\]
Let
\[
\partial:C\longrightarrow K^n,
\qquad
\partial(v)=(\beta_1(v),\dots,\beta_n(v)).
\]
Since \(\ker\partial=C\cap U_c\), rank--nullity and \(\dim C=n\) give
\[
\operatorname{rank}\partial+\dim(C\cap U_c)=n,
\qquad \operatorname{rank}\partial=k.
\]
Choose a set \(J\) of \(k\) block coordinates for which the coordinate
projection
\[
\operatorname{im}\partial\longrightarrow K^J
\]
is an isomorphism, and permute the blocks so that \(J=\{1,\dots,k\}\).
Choose pivot rows \(g_1,\dots,g_k\in C\) whose images under \(\partial\)
restrict to the standard basis of \(K^J\), and choose rows
\(S_1,\dots,S_r\) forming a basis of \(\ker\partial\).  These \(n\) rows form
a basis of \(C\): applying \(\partial\) annihilates the \(S_s\) and determines
all pivot coefficients, after which independence of the \(S_s\) determines
the remaining coefficients.

For \(i\ne j\), let \(b_{ij}\) be the first coordinate of the \(j\)-th
pivot block of \(g_i\), and let \(\ell_i\) be the terminal part of \(g_i\).
The terminal first coordinates of the \(S_s\) form \(D\).  Let \(d_{sj}\)
be the first coordinate of the
\(j\)-th pivot block of \(S_s\).  The block identity
\[
(\alpha,c\alpha+\beta)\cdot(\gamma,c\gamma+\delta)
=c(\alpha\delta+\beta\gamma)+\beta\delta
\]
and \(g_i\cdot S_s=0\) give
\[
d_{si}+\sum_{t=1}^r(\partial_c\ell_i)_tD_{st}=0.
\]
Thus \(d_{si}=-(D\partial_c\ell_i)_s\), which gives
\eqref{eq:universal-rank-box} without an additional coefficient matrix.
It remains to prove that \(D\) is nonsingular.  For \(S\in\ker\partial\), write
\(S=(a_1(1,c),\dots,a_n(1,c))\) and put \(a(S)=(a_1,\dots,a_n)\). Since
\[
(1,c)\cdot(s,t)=s+ct=c(t-cs),
\]
orthogonality of \(S\) to every \(v\in C\) gives
\[
0=S\cdot v=c\,a(S)\cdot\partial(v).
\]
Hence
\[
a(\ker\partial)\subseteq(\operatorname{im}\partial)^\perp.
\]
The map \(a\) is injective on \(\ker\partial\), and both spaces have dimension
\(r=n-k\); therefore equality holds.  If
\(a\in(\operatorname{im}\partial)^\perp\) vanishes outside \(J\), then the
surjectivity of \(\operatorname{im}\partial\to K^J\) forces \(a=0\).
Thus projection of \((\operatorname{im}\partial)^\perp\) onto the complementary
\(r\) coordinates is an isomorphism.  In the bases \(S_1,\dots,S_r\), its
matrix is precisely \(D\), proving \(\det D\ne0\).

The self-product of the \(i\)-th pivot row is
\(1+2ca_i+\ell_i\boldsymbol\cdot\ell_i=0\), so its diagonal block is forced
by \(\ell_i\).  For \(i\ne j\), pivot--pivot orthogonality is exactly
\eqref{eq:universal-rank-pp}; products between the \(S_s\) vanish
identically.  This proves the existence direction.  Conversely, the displayed
lower-left coefficients make pivot--terminal products vanish identically,
the definition of \(a_i\) gives the diagonal products, and
\eqref{eq:universal-rank-pp} gives the off-diagonal products.  Applying the
linear functional \((s,t)\mapsto t-cs\) blockwise to the first
\(k\) block columns reads off the pivot coefficients; after they vanish, the
remaining relation among the last \(r\) rows is killed by the invertibility of \(D\).
Thus all \(k+r\) rows are independent.  Their span is a totally isotropic
subspace of half the ambient dimension, hence is self-dual.

For the recursive assertion, restriction to \(I\) preserves every
\(\ell_i\), every \(a_i\), and each relation
\eqref{eq:universal-rank-pp}; the corresponding lower-left column remains
\(-D\partial_c\ell_i\).  Since \(D\) is unchanged, the
converse already proved applies.
For \(I=\varnothing\), the remaining rows are \((D_{st}(1,c))_{t=1}^r\).
As \(\operatorname{rowspan}D=K^r\), their span is exactly \(U_c^{(r)}\).
\end{proof}

Thus, writing \(G_{k,r}\) for the displayed matrix, each first pivot
can be removed in the same way:
\[
G_{k,r}=
\left[\begin{array}{c|c}
(a_1,ca_1+1) & \rho_1\\ \hline
\gamma_1(1,c) & G_{k-1,r}
\end{array}\right],
\qquad
G_{k,r}\longrightarrow G_{k-1,r}
\longrightarrow\cdots\longrightarrow G_{0,r}=(D_{st}(1,c)).
\]
Here \(\rho_1=(b_{12}(1,c),\ldots,b_{1k}(1,c),\ell_1)\) and
\(\gamma_1=(b_{21},\ldots,b_{k1},
-(D\partial_c\ell_1)_1,\ldots,-(D\partial_c\ell_1)_r)^{\mathsf T}\).
Each arrow deletes a pivot row together with its two-coordinate block
column; it is not puncturing the whole code.  Each retained row space is
self-dual.  This is the repeated boxed structure, with the same nonsingular
terminal block at every step.

The integer \(r\) records \(C\cap U_c\) for the chosen ordered coordinate
pairing.  It is intrinsic after the blocks and their orientations have been
fixed, but a different scalar-coordinate permutation can change it.
The entries of \(D\), however, depend on the chosen basis of \(C\cap U_c\).
Multiplying the last \(r\) rows by \(D^{-1}\) normalizes their terminal
coefficient matrix to \(I_r\) without changing the code.  We retain \(D\) only
to display the terminal basis carried unchanged through every repeated step.
Universality
therefore requires variable \(r\) and the unrestricted pivot and terminal
coefficients, not a nonidentity choice of \(D\).

\begin{lemma}[Building-up with fixed terminal matrix]
\label{lem:rank-boxed-building-up}
Let \(b,\ell,D\) satisfy the hypotheses of the converse in
Theorem~\ref{thm:universal-rank-boxed}; thus the parent is
\(G(c;b,\ell,D)\).  Choose a terminal row
\(\ell_0\in(K^2)^r\) and arbitrary coefficients
\(u_1,\ldots,u_k\in K\).  Put
\[
a_0=\frac c2(1+\ell_0\mathbin{\boldsymbol\cdot}\ell_0),
\qquad b^+_{0i}=u_i,
\qquad b^+_{i0}=c(\ell_i\mathbin{\boldsymbol\cdot}\ell_0)-u_i,
\]
and retain \(b^+_{ij}=b_{ij}\) for \(i,j\geq1\).  With terminal rows
\((\ell_0,\ell_1,\ldots,\ell_k)\) and the same \(D\), these coefficients
define a matrix \(G_{k+1,r}\) whose row space is self-dual.
Deleting its first row and first two-coordinate block column gives
\(G_{k,r}\) exactly. This construction may be iterated arbitrarily
many times, keeping \(r\) and \(D\) fixed.
\end{lemma}

\begin{proof}
Since \(K\) has odd characteristic, \(2\) is invertible and \(a_0\) is
well defined.  The new diagonal row is isotropic by its definition, and
for each old pivot row,
\[
c(b^+_{0i}+b^+_{i0})
 +\ell_0\mathbin{\boldsymbol\cdot}\ell_i
=c^2(\ell_i\mathbin{\boldsymbol\cdot}\ell_0)
 +\ell_0\mathbin{\boldsymbol\cdot}\ell_i=0.
\]
All old relations are unchanged.  The new lower-left column is
\(-D\partial_c\ell_0\), while every old column and the
terminal matrix \(D\) remain unchanged.  The converse in
Theorem~\ref{thm:universal-rank-boxed} therefore gives independence and
self-duality.  The resulting data satisfy the same hypotheses, so the
construction can be repeated. Finally, the defect
map on the \((k+1+r)\)-dimensional code has rank \(k+1\), so its kernel has
dimension~\(r\).
\end{proof}

The next two lemmas identify precisely how this repeated construction
corresponds to a Kim--Lee step with the same parent matrix. All changes
below are invertible row operations; the coordinate order and Euclidean
form remain fixed.

\begin{lemma}[Repeated steps and Kim--Lee building-up]
\label{lem:repeated-kim-lee}
Use the notation and hypotheses of Lemma~\ref{lem:rank-boxed-building-up},
and write \(G=G_{k,r}\), with rows \(g_1,\ldots,g_m\), where
\(m=k+r\). Put
\[
\gamma=\begin{pmatrix}b^+_{10}\\ \vdots\\ b^+_{k0}\\
        -D\partial_c\ell_0\end{pmatrix},\qquad
\rho=\bigl(u_1(1,c),\ldots,u_k(1,c),\ell_0\bigr).
\]
Then the successor is exactly
\begin{equation}\label{eq:repeated-border}
G^+=G_{k+1,r}=
\begin{pmatrix}
a_0&ca_0+1&\rho\\
\gamma_1&c\gamma_1&g_1\\
\vdots&\vdots&\vdots\\
\gamma_m&c\gamma_m&g_m
\end{pmatrix}.
\end{equation}
Its row space has a Kim--Lee presentation with this parent and these
first two coordinates if and only if \(\gamma\ne0\), equivalently,
\[
\neg\bigl(\partial_c\ell_0=0\ \text{and}\
u_i=c(\ell_i\mathbin{\boldsymbol\cdot}\ell_0)
\text{ for every }i\bigr).
\]
More precisely, if \(\gamma_s\ne0\), set
\begin{equation}\label{eq:repeated-kim-normalization}
\lambda=\frac{c-a_0}{\gamma_s},\qquad
x=c^{-1}(\rho+\lambda g_s).
\end{equation}
Replacing the first row \(b_0\) of \(G^+\) by
\(c^{-1}(b_0+\lambda b_s)\), and leaving all lower rows unchanged,
gives the exact matrix
\begin{equation}\label{eq:repeated-kim-matrix}
\begin{pmatrix}
1&0&x\\
-x\cdot g_1&-c(x\cdot g_1)&g_1\\
\vdots&\vdots&\vdots\\
-x\cdot g_m&-c(x\cdot g_m)&g_m
\end{pmatrix},\qquad x\cdot x=-1.
\end{equation}
If \(\gamma=0\), instead,
\[
a_0=c/2,\qquad \rho\in\operatorname{row}(G),\qquad
\operatorname{row}(G^+)=K(1,-c)\oplus\operatorname{row}(G)
\]
in the displayed coordinates. This last code is not the row space of
\eqref{eq:repeated-kim-matrix} for any \(x\) with this fixed parent and
coordinate order.
\end{lemma}

\begin{proof}
The coefficient formulas in Lemma~\ref{lem:rank-boxed-building-up}
give \eqref{eq:repeated-border} entry by entry. Since \(G\) and \(G^+\)
generate self-dual codes, their rows are pairwise orthogonal. Using
\(c^2=-1\), we obtain
\begin{equation}\label{eq:repeated-border-gram}
\rho\cdot g_i=-c\gamma_i,\qquad
\rho\cdot\rho=-1-2ca_0.
\end{equation}
Indeed, the first two coordinates contribute \(c\gamma_i\) to
\(b_0\cdot b_i\), and \(1+2cp\) to \(b_0\cdot b_0\).
If \(\gamma_s\ne0\), then
\[
a_0+\lambda\gamma_s=c,\qquad
ca_0+1+\lambda c\gamma_s=c^2+1=0.
\]
Thus the new first pair is \((1,0)\). Moreover,
\[
x\cdot g_i
=c^{-1}(\rho\cdot g_i+\lambda g_s\cdot g_i)
=-\gamma_i.
\]
The new first row is isotropic, being a linear combination of orthogonal
rows, so \(1+x\cdot x=0\). This proves the literal equality
\eqref{eq:repeated-kim-matrix}. The row operation has determinant
\(c^{-1}\ne0\), hence preserves the row space, and every parent tail
is unchanged. In the sign convention of
Theorem~\ref{thm-q-ary-build-1}, the square root used in this Kim--Lee
matrix is \(-c\).

If \(\gamma=0\), \eqref{eq:repeated-border-gram} gives
\(\rho\in\operatorname{row}(G)^\perp=\operatorname{row}(G)\).
Consequently \(\rho\cdot\rho=0\), and \(1+2ca_0=0\) gives \(a_0=c/2\).
Subtracting a combination of the lower rows removes \(\rho\) from
\(b_0\). Its remaining first pair is
\((c/2,1/2)=(c/2)(1,-c)\), proving the direct-sum equality.
Every vector in this direct sum has its first pair in \(K(1,-c)\),
which does not contain \((1,0)\). The first row of
\eqref{eq:repeated-kim-matrix} therefore cannot belong to this code.
Finally, \(D\) is invertible, so
\(D\partial_c\ell_0=0\) forces
\(\partial_c\ell_0=0\).  The remaining entries of \(\gamma\) vanish
exactly when
\(u_i=c(\ell_i\mathbin{\boldsymbol\cdot}\ell_0)\) for every \(i\).
\end{proof}

\begin{lemma}[Kim--Lee steps in repeated form]
\label{lem:kim-lee-repeated}
Fix a parent \(G=G_{k,r}\) satisfying the hypotheses of
Lemma~\ref{lem:rank-boxed-building-up}. Every \(x\in K^{2(k+r)}\)
with \(x\cdot x=-1\) gives a Kim--Lee matrix
\eqref{eq:repeated-kim-matrix} that can be changed, by an invertible
operation on its first row alone, into a successor from
Lemma~\ref{lem:rank-boxed-building-up}. The parent matrix and \(D\)
are unchanged, and the resulting \(\gamma\) is nonzero.
\end{lemma}

\begin{proof}
Write \(\gamma_i=-x\cdot g_i\) and
\(b_i=(\gamma_i,c\gamma_i,g_i)\). For the old pivot blocks put
\[
a_j=c(x_{2j}-cx_{2j-1})\quad(1\le j\le k),\qquad
b'_0=c(1,0,x)-\sum_{j=1}^k a_j b_j.
\]
Its tail and first pair are
\[
\rho=cx-\sum_{j=1}^k a_j g_j,\qquad
(a_0,ca_0+1),\qquad a_0=c-\sum_{j=1}^k a_j\gamma_j.
\]
The old pivot blocks satisfy
\((g_i)_{2j}-c(g_i)_{2j-1}=\delta_{ij}\) for \(1\le i,j\le k\).
Hence
\[
\rho_{2j}-c\rho_{2j-1}
=c(x_{2j}-cx_{2j-1})-a_j=0\quad(1\le j\le k).
\]
Set \(u_j=\rho_{2j-1}\), and let \(\ell_0\) be the terminal part of
\(\rho\). Orthogonality of \(b'_0\) to the terminal rows and then to the
pivot rows gives, respectively,
\[
 (\gamma_{k+1},\ldots,\gamma_{k+r})^{\mathsf T}
 =-D\partial_c\ell_0,\qquad
 \gamma_i=c(\ell_i\mathbin{\boldsymbol\cdot}\ell_0)-u_i.
\]
Finally, isotropy of \(b'_0\) gives
\[
0=1+2ca_0+\ell_0\mathbin{\boldsymbol\cdot}\ell_0,
\qquad
a_0=\frac c2\bigl(1+\ell_0\mathbin{\boldsymbol\cdot}\ell_0\bigr).
\]
These are exactly the successor coefficients of
Lemma~\ref{lem:rank-boxed-building-up}. No lower row has changed, and
the first-row operation has determinant \(c\ne0\).
If \(\gamma=0\), then \(x\) is orthogonal to every row of \(G\),
so self-duality gives \(x\in\operatorname{row}(G)\) and
\(x\cdot x=0\), contradicting \(x\cdot x=-1\).
\end{proof}

Thus, with the parent and coordinate pair fixed, repeated extensions are
exactly Kim--Lee steps together with the length-two direct-sum additions
identified above. Different choices of \(s\) in
\eqref{eq:repeated-kim-normalization} need not give the same vector \(x\);
the statements concern generator presentations, not a bijection of
parameter tuples. The direct-sum exception has minimum distance at most
two, and so is excluded when the child has \(d>2\). A different coordinate
permutation can change the parent reduction; no obstruction to such a
different presentation is asserted.

\begin{lstlisting}[caption={Exact Lean goals for Theorem~\ref{thm:universal-rank-boxed}.}]
theorem paperRankBoxedRows_forward_selfDual {k r : Nat} (c : K)
    (b : Fin k -> Fin k -> K)
    (ell : Fin k -> Fin r -> SplitBlock K)
    (D : Fin r -> Fin r -> K)
    (hc : c * c = -1) (h2 : (2 : K) ≠ 0)
    (hD : RankBoxCoreFullRank D)
    (hoff : PaperOffDiagonalRelations c b ell) :
    RankBoxedPairwiseOrthogonal (paperRankBoxedRows c b ell D) ∧
      LinearIndependent K (paperRankBoxedRows c b ell D) ∧
      rankBoxedRowSpace (paperRankBoxedRows c b ell D) =
        (rankBoxRowBilin (K := K) (k := k) (r := r)).orthogonal
          (rankBoxedRowSpace (paperRankBoxedRows c b ell D))

theorem every_qary_selfDualCode_has_rankBoxed_normalForm
    {n : Nat} (c : K) (hc : c ^ 2 = (-1 : K))
    (h2 : (2 : K) ≠ 0)
    {C : Submodule K (QaryBlockRow K (Fin n))}
    (hC : QaryBlockSelfDualCode C) :
    HasQaryRankBoxedNormalForm c C
\end{lstlisting}
In the listing, the first theorem is the forward direction for the displayed
matrix, and its hypothesis \texttt{hoff} is precisely
\eqref{eq:universal-rank-pp}.  The separate formal statements for repeated
steps prove the explicit normalization \eqref{eq:repeated-kim-normalization},
the direct-sum classification, and the reverse first-row operation in
Lemma~\ref{lem:kim-lee-repeated}.

\begin{example}[Rank-two building-up over \(\GF{5}\)]
\label{ex:rank-two-gf5}
For the oriented pairing displayed below, this example lies outside the
rank-one specialization of Corollary~\ref{thm:split-boxed-form}: its terminal
matrix has order two and its pivot diagonal blocks are not all \(01\).  It
illustrates the nonsingular matrix \(D\) and the forced diagonal terms in the
general rank-two coordinates.  It is not intended as a
permutation-invariant obstruction to a rank-one presentation.
Take \(c=2\), so \(c^2=-1\), and let \(C\) be generated by
\[
G_0=\begin{pmatrix}
1&3&3&2&1&1&1&3\\
3&4&2&1&4&4&3&2\\
1&4&3&4&2&3&4&2\\
1&2&2&2&3&4&4&4
\end{pmatrix}.
\]
All calculations in this example are in \(\GF{5}\).
One checks that \(G_0G_0^{\mathsf T}=0\), while the minor formed by the first
four columns has determinant \(4\). Hence \(C\) has dimension \(4\) and is
self-dual. None of the entries of \(G_0\) is zero, and the repeated block
structure is not visible in this presentation.

Pair adjacent coordinates and let \(G_0^\sigma\) be obtained by ordering the
four two-coordinate blocks as \((2,4,1,3)\). The invertible row transformation
\[
T=\begin{pmatrix}
4&1&0&0\\
2&1&2&0\\
1&3&0&1\\
4&0&3&4
\end{pmatrix},\qquad \det T=3,
\]
gives the following matrix, written in two-coordinate blocks:
\[
G_{2,2}=TG_0^\sigma=
\left[\begin{array}{cc|cc}
(4,4)&(2,4)&(2,1)&(3,3)\\
(4,3)&(3,2)&(2,3)&(0,2)\\ \hline
(1,2)&(4,3)&(1,2)&(1,2)\\
(4,3)&(2,4)&(1,2)&(2,4)
\end{array}\right].
\]
This is \eqref{eq:universal-rank-box} with
\[
\ell_1=((2,1),(3,3)),\qquad
\ell_2=((2,3),(0,2)),\qquad
b_{12}=2,\quad b_{21}=4,
\]
and
\[
D=\begin{pmatrix}1&1\\1&2\end{pmatrix},\qquad \det D=1.
\]
Indeed, the forced diagonal coefficients are
\(a_1=4\) and \(a_2=3\), giving the blocks \((4,4)\) and \((3,2)\).
Moreover,
\(-D\partial_2\ell_1=(1,4)^{\mathsf T}\) and
\(-D\partial_2\ell_2=(4,2)^{\mathsf T}\), exactly as in
the lower-left block columns.  The sole off-diagonal relation
\eqref{eq:universal-rank-pp} is also satisfied.

Applying \((s,t)\mapsto t-2s\) to each block of \(G_{2,2}\) gives
\[
\partial G_{2,2}=
\begin{pmatrix}
1&0&2&2\\
0&1&4&2\\
0&0&0&0\\
0&0&0&0
\end{pmatrix}.
\]
Its rank is \(2\), so \(r=\dim(C\cap U_2)=4-2=2\).
Here the block permutation preserves \(U_2\), so the same value holds in
the original coordinates. Consequently, whole-block permutations and row
operations alone cannot turn this example into a rank-one form for this
chosen oriented pairing.

The qualifier about the pairing is essential.  Reverse only the first
adjacent block, so that the ordered pairs are
\((2,1),(3,4),(5,6),(7,8)\).  For the resulting coordinate permutation
\(\tau\), exact row reduction gives
\[
\partial G_0^\tau\ \sim\
\begin{pmatrix}
1&0&0&0\\
0&1&0&4\\
0&0&1&3\\
0&0&0&0
\end{pmatrix}.
\]
Thus the defect rank is \(3\), so \(r=1\) for this new pairing.  In
particular, this same code is permutation equivalent to a rank-one boxed
representative.  Here ``chosen pairing'' means that both the partition into
two-coordinate blocks and the orientation inside every block have already
been specified.

Now delete the first pivot row together with the first block column, and
repeat once more. The retained matrices are
\[
G_{1,2}=
\left[\begin{array}{c|cc}
(3,2)&(2,3)&(0,2)\\ \hline
(4,3)&(1,2)&(1,2)\\
(2,4)&(1,2)&(2,4)
\end{array}\right],\qquad
G_{0,2}=
\begin{pmatrix}
(1,2)&(1,2)\\
(1,2)&(2,4)
\end{pmatrix}.
\]
Theorem~\ref{thm:universal-rank-boxed} gives a sequence of Euclidean
self-dual row spaces with parameters
\[
[8,4]\longrightarrow[6,3]\longrightarrow[4,2].
\]
The same terminal matrix \(D\) remains throughout, and
\(\operatorname{rowspan}G_{0,2}=K(1,2)\oplus K(1,2)\).
Each step retains only the remaining generator rows before deleting their
first two coordinates; it is not puncturing the whole preceding code.

Starting instead from the same \([8,4]\) matrix \(G_{2,2}\),
we can build up twice to length \(12\). In Lemma~\ref{lem:rank-boxed-building-up}, choose
\[
\begin{array}{c|cc}
\text{step}&\ell_0&u\\ \hline
G_{2,2}\longrightarrow G_{3,2}&((0,0),(0,2))&(0,1)\\
G_{3,2}\longrightarrow G_{4,2}&((0,0),(0,2))&(0,0,2)
\end{array}
\]
with each new pivot prepended. Both steps give \(a_0=0\), so each newly
adjoined diagonal block is \(01\).  The final matrix is
\[
G_{4,2}=
\left[\begin{array}{cc|cc|cc}
(0,1)&(0,0)&(0,0)&(2,4)&(0,0)&(0,2)\\
(3,1)&(0,1)&(0,0)&(1,2)&(0,0)&(0,2)\\ \hline
(2,4)&(2,4)&(4,4)&(2,4)&(2,1)&(3,3)\\
(1,2)&(2,4)&(4,3)&(3,2)&(2,3)&(0,2)\\ \hline
(3,1)&(3,1)&(1,2)&(4,3)&(1,2)&(1,2)\\
(1,2)&(1,2)&(4,3)&(2,4)&(1,2)&(2,4)
\end{array}\right].
\]
For this displayed pairing, the largest repeated matrix is genuinely the
rank-two, rather than the rank-one, form.  Indeed,
\[
\partial G_{4,2}=
\begin{pmatrix}
1&0&0&0&0&2\\
0&1&0&0&0&2\\
0&0&1&0&2&2\\
0&0&0&1&4&2\\
0&0&0&0&0&0\\
0&0&0&0&0&0
\end{pmatrix},
\qquad
\operatorname{rank}(\partial G_{4,2})=4,
\]
and hence \(r=6-4=2\).  Thus the largest member of the displayed repeated
chain is not a rank-one box in these ordered coordinates.
The first two rows and block columns are the new ones. Their removal
leaves exactly the earlier \(G_{2,2}=TG_0^\sigma\), not a different
\([8,4]\) example. Removing only the first row and first block column
gives \(G_{3,2}\).  Hence \(G_{3,2}\) has one new \(01\) diagonal block and
\(G_{4,2}\) has two.  This is maximal while the original \(G_{2,2}\), whose
two inherited pivot diagonal blocks are \((4,4)\) and \((3,2)\), is retained
literally.  Swapping the coordinates of each newly adjoined \(01\) block
recovers the literal Kim--Lee heads \(10\) with slope \(3=-2\), as checked in
Lean. Thus, with arrows denoting building-up and its
inverse row-and-block reduction, we have
\[
G_{2,2}\ ([8,4])\ \longleftrightarrow\
G_{3,2}\ ([10,5])\ \longleftrightarrow\
G_{4,2}\ ([12,6]).
\]
Their exact parameters are \([8,4,4]\), \([10,5,4]\), and \([12,6,4]\).
All three row spaces are Euclidean self-dual, and all have \(r=2\) and the
same terminal matrix \(D\).  Thus this is a structural example rather than a
distance search; Section~\ref{sec:applications} subsequently applies the same
reduction--reconstruction mechanism to rank-one presentations of optimal
codes.

These two extensions also give the literal Kim--Lee construction, not just
matrices satisfying the Gram equations. Reverse the two coordinates in the
newly prepended \(01\) block; its first row then begins with \(10\), while the
parent coordinates remain unchanged. The corresponding tail vectors are
\[
x_1=(0,0,1,2,0,0,0,2),\qquad
x_2=(0,0,0,0,2,4,0,0,0,2),
\]
and \(x_1\cdot x_1=x_2\cdot x_2=4=-1\).
For each parent row \(g_i\), the reordered new head is
\((-y_i,-3y_i)\), where \(y_i=x_j\cdot g_i\). Thus it is exactly the
internal \(\operatorname{buildRows}(x_j,3,G)\) convention, with
\(3=-2\); the displayed boxed matrices themselves use \(c=2\).
This is a scalar-coordinate permutation, not a change of the Euclidean form.
Lemma~\ref{lem:rank-boxed-building-up}, beyond these two numerical steps,
ensures further extensions with the same \(D\), and
Lemma~\ref{lem:kim-lee-repeated} supplies the converse for every Kim--Lee
extension of either parent.
\end{example}

\subsection{Binary rank-one refinement}
\label{sec:binary-refinement}

The coefficient \(c/2\) in the preceding theorem is specific to odd
characteristic.  In characteristic two,
Theorem~\ref{thm:binary-universal-rank-one} gives the sharper statement
directly: after a scalar-coordinate permutation, \(r=1\), \(D=(1)\), every
pivot diagonal block is \(01\), every terminal pivot block is \(10\), and
the forced lower-left coefficient is (1).  Thus the result is the literal
binary box, rather than a formal substitution into the odd-characteristic
diagonal formula.
This permutation may change the original pairing of coordinates; it does
not assert \(r=1\) for every fixed pairing.

The normalization in \eqref{eq:repeated-kim-normalization} does not use
division by two. In a normalized binary CZ step, \(c=1\), the diagonal
coefficient is zero,
and the all-ones last row has \(\gamma_s=1\). Thus
\(\lambda=1\), \(x=\rho+\mathbf1\), and the first-row operation is
exactly \(b_0\mapsto b_0+b_s\). It gives \(x\cdot x=1\) and
the binary Kim matrix, with the same parent tails. The nonzero last
entry of \(\gamma\) excludes the direct-sum exception at every such
step. This specialization uses the first-row normalization, not the
odd-characteristic coefficient formula involving \(c/2\).

In 2012, Chinburg and Zhang~\cite{ChinburgZhang} proved that up to permutation equivalence, all binary self-dual codes of length at least 4 arise from Hilbert pairings on rings of $S$-integers of $\mathbb Q$.

\begin{theorem}(\cite[Theorem 1.5]{ChinburgZhang})  \label{thm-binary-hilbert-2}
Up to permutation equivalence, all binary self-dual codes of length at least 4 arise from Hilbert pairings on rings of $S$-integers of $K= \mathbb Q$. In fact, each such code arises up to permutation equivalence from infinitely many different subsets $S$ of the places of $K=\mathbb Q$.
\end{theorem}

Here ``equivalence'' is used exactly in the sense defined by
Chinburg--Zhang: a bijection of the chosen orthonormal bases whose linear
extension carries one code to the other. Thus it is coordinate-permutation
equivalence, not arbitrary monomial equivalence.

In the course of their proof, Chinburg--Zhang exhibit a boxed normal form for the resulting binary self-dual code. In this form, the bottom row is the all-ones vector, the distinguished terminal block-column is of type \(10\) above the bottom row, the diagonal block-entries are of type \(01\), and every remaining off-diagonal block is either \(00\) or \(11\), paired so that the two opposite blocks sum to \(11\).

In the table below, \(W_v=\mathbb Q_v^\times/(\mathbb Q_v^\times)^2\)
is the local square-class space with its Hilbert pairing. The primes
\(p_i\) are distinct and congruent to \(3\pmod4\); the column headings
list local bases and the rows list global \(S\)-units~\cite{ChinburgZhang}.

\begin{center}
\begin{tabular}{|c|ccccc|}
\toprule
$S$-units $\backslash$  places & $W_{p_1}$  & $W_{p_2}$ & $\cdots$ & $W_2$  & $W_{\mathbb R}$ \\
 & $\{-p_1, p_1\}$ & $\{-p_2, p_2\}$ & $\cdots$  & $\{-2, -10, -5 \}$ & $\{-1 \}$  \\
\midrule
$p_1$  & 01   & 00/11  &  & 00/11 ~~1  &  0 \\
$p_2$  & 11/00   & 01  &  & 00/11 ~~1  &  0 \\
$\vdots$ & $\vdots$ &   &  $\ddots$ &   & $\vdots$ \\
2  & 11/00   & 11/00  &   & 01 ~~~~~~~1  &  0 \\
$-1$ & 11      & 11 & $\cdots$ & 11 ~~~~~~ 1 &  1 \\
\bottomrule
\end{tabular}
\end{center}

The point of this boxed form is twofold: every binary self-dual code of length \(2n\) admits such a representative, and deleting the distinguished two-coordinate block together with the top row produces a shorter self-dual code of length \(2n-2\). This is a coordinate reduction in the standard Euclidean form. The deleted coordinate plane has Gram matrix \(I_2\); it is not an alternating hyperbolic plane.

The two-coordinate reduction itself does not require the entire matrix to have
already been put in boxed form. It retains the codewords whose first two
coordinates are equal and then deletes those coordinates. This differs from
ordinary coordinate shortening, which retains only codewords that are zero
at the deleted positions. The following lemma gives the reduction used in
the induction.

\begin{samepage}
\begin{lemma}[Binary two-coordinate reduction]
\label{lem:binary-two-coordinate-reduction}
Let \(C=C^\perp\subseteq\mathbb F_2^{\,2+s}\), and suppose that
\(x\in C\) satisfies \(x_1=0\) and \(x_2=1\).  Put
\[
C_=\coloneqq\{v\in C:v_1+v_2=0\},
\qquad
\pi(v_1,v_2,v')=v',
\qquad
D\coloneqq\pi(C_=)\subseteq\mathbb F_2^s.
\]
Then \(\pi|_{C_=}\) is injective and \(D=D^\perp\).  In particular, \(s\)
is even and \(\dim D=s/2\).
\end{lemma}
\end{samepage}

\begin{proof}
If \(v\in C_=\) and \(\pi(v)=0\), then \(v=(a,a,0)\).  Since \(v,x\in C\)
and \(C\subseteq C^\perp\),
\[
0=x\mathbin{\cdot}v=0\cdot a+1\cdot a=a.
\]
Thus \(v=0\), proving injectivity.  If \(u=(a,a,g)\) and \(v=(b,b,h)\)
belong to \(C_=\), then
\[
0=u\mathbin{\cdot}v=2ab+g\mathbin{\cdot}h
  =g\mathbin{\cdot}h,
\]
so \(D\subseteq D^\perp\).  Finally, the functional
\(\delta:C\to\mathbb F_2\), \(\delta(v)=v_1+v_2\), is surjective because
\(\delta(x)=1\).  Hence
\[
\dim C_=\dim C-1=\frac{2+s}{2}-1=\frac{s}{2}.
\]
Injectivity of \(\pi|_{C_=}\) gives \(\dim D=s/2\); together with
\(D\subseteq D^\perp\), this yields \(D=D^\perp\).
\end{proof}

In the following listing, the restriction \(\pi|_{C_=}\) is denoted by
\texttt{binaryTwoCoordinateMap}. Its image \(D\) is denoted by
\texttt{binaryReducedCode}. The conclusions state injectivity,
self-duality, the dimension formula, and the evenness of the length.

\noindent\begin{minipage}{\linewidth}
\begin{lstlisting}[caption={Lean formulation of binary two-coordinate reduction.}]
theorem paper_binary_twoCoordinateReduction_exact {n : ℕ}
    {C : Submodule (ZMod 2) (Fin (2 + n) → ZMod 2)}
    {x : Fin (2 + n) → ZMod 2}
    (hC : paperSelfDualCode (K := ZMod 2) C)
    (hxC : x ∈ C) (hx0 : x 0 = 0) (hx1 : x 1 = 1) :
    Function.Injective (binaryTwoCoordinateMap C) ∧
      paperSelfDualCode (K := ZMod 2) (binaryReducedCode C) ∧
      2 * Module.finrank (ZMod 2) (binaryReducedCode C) = n ∧ Even n
\end{lstlisting}
\end{minipage}

\begin{theorem}[Universal binary rank-one boxed normal form]
\label{thm:binary-universal-rank-one}
For every \(N\geq 1\) and every Euclidean self-dual code
\(C\subseteq\mathbb F_2^{2N}\), there are a permutation \(\sigma\) of the
\(2N\) scalar coordinates and a matrix
\(b=(b_{ij})\in M_{N-1}(\mathbb F_2)\) such that
\begin{equation}
b_{ii}=0,
\qquad
b_{ij}+b_{ji}=1\quad(i\ne j).
\label{eq:binary-rank-one-laws}
\end{equation}
and \(\sigma(C)=\operatorname{rowspan}G(b)\), where, writing
\(a(11)=(a,a)\),
\begin{equation}
G(b)=
\left[
\begin{array}{ccccc|c}
01 & b_{12}(11) & b_{13}(11) & \cdots & b_{1,N-1}(11) & 10\\
b_{21}(11) & 01 & b_{23}(11) & \cdots & b_{2,N-1}(11) & 10\\
\vdots & \vdots & \ddots & & \vdots & \vdots\\
b_{N-1,1}(11) & b_{N-1,2}(11) & \cdots & & 01 & 10\\
\hline
11 & 11 & \cdots & & 11 & 11
\end{array}
\right].
\label{eq:binary-rank-one-box}
\end{equation}
For \(N=1\), this means the single row \(11\).
\end{theorem}

\begin{proof}
We induct on \(N\).  If \(N=1\), self-duality forces \(C\) to be the
unique isotropic line \(\langle(1,1)\rangle\subseteq\mathbb F_2^2\).
Assume \(N\geq2\).  For every \(v\in C\),
\[
v\mathbin{\cdot}{\bf1}=v\mathbin{\cdot}v=0;
\]
hence \({\bf1}\in C^\perp=C\).  Since \(\dim C=N>1\), some word of \(C\)
is nonconstant.  After a scalar-coordinate permutation it has the form
\(x=(0,1,x')\).  The binary two-coordinate reduction lemma gives the
self-dual predecessor
\[
D=\pi\{v\in C:v_1+v_2=0\}
   \subseteq\mathbb F_2^{2N-2}.
\]
By induction, after a permutation of the tail coordinates, \(D\) is generated
by a box \(G(b)\) with pivot rows \(g_1,\ldots,g_{N-2}\) and all-ones row
\(M={\bf1}\), satisfying \eqref{eq:binary-rank-one-laws}.

Put \(p=x+{\bf1}=(1,0,Z)\).  Since \(p\in C=C^\perp\),
\begin{equation}
0=p\mathbin{\cdot}p=1+Z\mathbin{\cdot}Z,
\qquad Z\mathbin{\cdot}Z=1.
\label{eq:binary-extension-norm}
\end{equation}
Exact reconstruction from the two-coordinate reduction gives
\begin{equation}
C=\operatorname{span}\left\{
A=(1,0,Z),\quad
\widehat g_i=(Z\mathbin{\cdot}g_i,Z\mathbin{\cdot}g_i,g_i),\quad
\widehat M=(Z\mathbin{\cdot}M,Z\mathbin{\cdot}M,M)
\right\}.
\label{eq:binary-kim-reconstruction}
\end{equation}
Thus this is literally the Kim extension of the inductively normalized
predecessor.

For an old two-coordinate pivot block \(j\), let
\(\delta_j(Z)=Z_{j,1}+Z_{j,2}\), and set
\begin{equation}
U=A+\sum_{j=1}^{N-2}\delta_j(Z)\widehat g_j,
\qquad
T=U+U_{\mathrm{term},2}\widehat M.
\label{eq:binary-row-correction}
\end{equation}
In the old box, the linear functional \((u,v)\mapsto u+v\) is one on the
diagonal block of \(g_j\), one on its terminal block, and zero on every
other pivot block and on \(M\).  Consequently every old pivot block of
\(T\) has coordinate sum zero and is therefore either \(00\) or \(11\).
Equation~\eqref{eq:binary-extension-norm} gives
\begin{equation}
T_{\mathrm{term}}=10.
\label{eq:binary-corrected-terminal}
\end{equation}
If \(s=T_{1,2}\), then the new diagonal block of \(T\) is
\((1+s,s)\).  It is already \(01\) when \(s=1\), and swapping the two new
scalar coordinates makes it \(01\) when \(s=0\).  This swap fixes the equal
new head blocks of every \(\widehat g_i\) and \(\widehat M\).  The latter is
the all-ones row, while all old diagonal and terminal blocks remain \(01\)
and \(10\), respectively.  Hence these corrected rows have exactly the form
\eqref{eq:binary-rank-one-box} for a matrix \(b'\).

Row operations and the optional pair swap preserve the code and its Euclidean
inner product.  For two distinct pivot rows of the resulting box, direct block
multiplication gives
\[
0=g_i(b')\mathbin{\cdot}g_j(b')
  =b'_{ij}+b'_{ji}+1,
\]
so \(b'_{ij}+b'_{ji}=1\); the diagonal blocks give \(b'_{ii}=0\).
Composing the initial pivot permutation, the inductive tail permutation, and
the optional new-pair swap yields the required single scalar-coordinate
permutation \(\sigma\).
\end{proof}

\noindent\begin{minipage}{\linewidth}
\begin{lstlisting}[caption={Lean statement for the universal binary rank-one form.}]
theorem binarySelfDualCode_has_rankOneNormalForm
    {k : Nat}
    {C : Submodule (ZMod 2) (Fin (2 * (k + 1)) → ZMod 2)}
    (hC : paperSelfDualCode C) :
    HasBinaryCzRankOneNormalForm C
\end{lstlisting}
\end{minipage}

The similarity between Theorem~\ref{thm-binary-build-1} and this boxed matrix is not accidental. Kim's theorem is a bottom-up extension statement, while the Chinburg--Zhang boxed form is a top-down reduction statement. The next theorem shows that the extension and reduction are inverse constructions.

For the arithmetic comparison, let \(K\) be a number field, let \(S\)
be a finite set containing its archimedean and dyadic places, and put
\(U=\operatorname{Spec}\mathcal O_{K,S}\). Set
\[
W=\bigoplus_{v\in S}H^1_{\mathrm{et}}(K_v,\mu_2),\qquad
\Phi:H^1_{\mathrm{et}}(U,\mu_2)\longrightarrow W,\qquad L=\operatorname{im}\Phi.
\]
At real places the local groups use Tate cohomology. The pairing
\(B_{\rm coh}\) is the sum of the local Hilbert pairings, equivalently the
cup product followed by the compact-support boundary map
\cite[Theorem~1.1]{ChinburgZhang}; it is nondegenerate and
\(L=L^{\perp_{B_{\rm coh}}}\).
If \(-1\) is nonsquare in every non-complex \(K_v\) for \(v\in S\),
local orthonormal bases give an ordered basis \(E\) of \(W\)
\cite[Theorem~1.2 and Corollary~1.4]{ChinburgZhang}.
Writing \(\dim W=2n\), its coordinate map is a form isometry
\[
\iota_E:(W,B_{\rm coh})\xrightarrow{\sim}
(\mathbb F_2^{2n},\langle\ ,\ \rangle_{\rm E}).
\]
The following comparison uses only these linear-algebraic data.

\begin{theorem}[Euclidean realization and binary boxed reduction]\label{thm:binary-cz-kim}
Let $n\ge 2$, and let $L=\operatorname{im}\Phi\subset W$ be the Chinburg--Zhang arithmetic image and suppose
\[
L=L^{\perp_{B_{\rm coh}}}.
\]
Fix the form isometry $\iota_E$ above and put $C=\iota_E(L)$. Apply
Theorem~\ref{thm:binary-universal-rank-one} and add the all-ones row to
the first pivot row. The resulting permutation-equivalent representative
of $C$ has generator matrix
\[
G'=
\begin{pmatrix}
1&0&x\\
y_1&y_1&g_1\\
\vdots&\vdots&\vdots\\
y_{n-1}&y_{n-1}&g_{n-1}
\end{pmatrix},
\qquad
x,g_i\in\mathbb F_2^{2n-2}.
\]
Let $G$ be the matrix with rows $g_1,\ldots,g_{n-1}$ and let $C'=\operatorname{rowspan}(G)$.  Then:
\begin{enumerate}
\item $L$ is self-orthogonal for $B_{\rm coh}$ and $C$ is Euclidean self-dual.
\item In the displayed Euclidean coordinates,
\[
\mathbb F_2^{2n}=E_2\perp\mathbb F_2^{2n-2},
\qquad \operatorname{Gram}(E_2)=I_2.
\]
The plane $E_2$ is non-alternating and is not isometric to the alternating hyperbolic plane.  Nevertheless, $\langle(1,1)\rangle\subset E_2$ is an isotropic line and $\langle(1,0),(1,1)\rangle_{\rm E}=1$.
\item The rows $g_i$ are linearly independent and pairwise orthogonal; hence $C'$ is a Euclidean self-dual code of length $2n-2$.
\item For every $i$,
\[
y_i=\langle x,g_i\rangle_{\rm E}.
\]
Consequently $G'$ is exactly the Kim building-up matrix obtained from $G$ and $x$.
\item Deleting the top row and the first two coordinates of $G'$ gives $G$, and rebuilding from $G$ gives $G'$ exactly.  Therefore, after restoring the coordinate permutation used to obtain the boxed representative, the Chinburg--Zhang boxed reduction and the Kim extension are inverse up to permutation equivalence.
\end{enumerate}
\end{theorem}

\begin{proof}
The assumed equality \(L=L^{\perp_{B_{\rm coh}}}\) implies self-orthogonality.
Since \(\iota_E\) is a form isometry, it carries orthogonal complements to
orthogonal complements, so \(C=\iota_E(L)\) is Euclidean self-dual.
Theorem~\ref{thm:binary-universal-rank-one} gives a boxed representative
whose first pivot row starts with \(01\) and whose all-ones row starts
with \(11\). Adding the latter to the former changes that head to
\(10\). Every other row has equal first coordinates. This invertible
row operation gives precisely the displayed matrix \(G'\), without any
further change of coordinates. Composing \(\iota_E\) with the selected
coordinate permutation gives a form isometry whose image of \(L\) is
exactly the displayed row space.

The first two standard coordinate vectors have Gram matrix $I_2$.  Thus $E_2$ is non-alternating because $(1,0)$ has self-product $1$.  The alternating hyperbolic plane has Gram matrix $\left(\begin{smallmatrix}0&1\\1&0\end{smallmatrix}\right)$, and alternation is invariant under form isometry, so these planes are not isometric.  On the other hand,
\[
(1,1)\mathbin{\cdot}(1,1)=0,
\qquad
(1,0)\mathbin{\cdot}(1,1)=1,
\]
which gives the stated isotropic line in \(E_2\) without turning $E_2$ into a hyperbolic plane.

Self-duality of $C$ gives orthogonality of the rows of $G'$. In particular,
\(0=(1,0,x)\cdot(1,0,x)=1+x\cdot x\), so \(x\cdot x=1\).
Moreover,
\[
0=(1,0,x)\mathbin{\cdot}(y_i,y_i,g_i)
  =y_i+x\mathbin{\cdot}g_i,
\]
and characteristic two yields $y_i=x\mathbin{\cdot}g_i$. Therefore $G'$ is precisely the generator family in Kim's building-up construction. Since $C$ has dimension $n$ and $G'$ has $n$ rows, those rows are linearly independent. The linear map
\[
g\longmapsto(x\mathbin{\cdot}g,\ x\mathbin{\cdot}g,\ g)
\]
has tail projection as a left inverse, so the rows of $G$ are linearly independent.  For all $i,j$,
\[
(y_i,y_i,g_i)\mathbin{\cdot}(y_j,y_j,g_j)
=2y_iy_j+g_i\mathbin{\cdot}g_j
=g_i\mathbin{\cdot}g_j.
\]
Thus the $g_i$ are pairwise orthogonal.  Their span has dimension $n-1$, half the ambient length $2n-2$, and is therefore self-dual.  Finally, tail projection gives $G$ after deleting the top row and first coordinate pair, while the coefficient identity rebuilds $G'$ exactly.  The only remaining change from the original code is the coordinate permutation used to reach the boxed representative.
\end{proof}



\noindent\begin{minipage}{\linewidth}
In the following listing, \texttt{e} includes the selected permutation of
the orthonormal coordinates. Thus \texttt{hcode} is equality with the
actual image of \(L\). The internal predicate \texttt{IsSelfOrthogonal}
denotes equality with the orthogonal complement, as required by the
hypothesis on \(L\).

\begin{lstlisting}[caption={Lean for Theorem~\ref{thm:binary-cz-kim}.}]
theorem dot_r0_riBin_eq_zero_iff [CharP K 2] :
    dot (r0 x) (riBin y g) = 0 <-> y = dot x g

theorem boxedFamily_tail_paperSelfDualCode [CharP K 2]
    (hself : paperSelfDualCode (rowSpace (boxedFamily x Y G))) :
    paperSelfDualCode (rowSpace G)

theorem paper_binary_cz_kim_corrected [CharP K 2]
    (he : IsFormIsometry Bcoh dotBilin e)
    (hL : IsSelfOrthogonal Bcoh L)
    (hcode : L.map e.toLinearMap =
      rowSpace (boxedFamily x Y G)) :
    IsSelfOrthogonal dotBilin (L.map e.toLinearMap) ∧
    paperSelfDualCode (rowSpace G) ∧
    (∀ i, Y i = dot x (G i)) ∧
    boxedFamily x Y G = buildRowsBin x G ∧
    deleteBinaryHeadPair (boxedFamily x Y G) = G ∧
    CodeEquiv
      (buildRowsBin x (deleteBinaryHeadPair (boxedFamily x Y G)))
      (boxedFamily x Y G)
\end{lstlisting}
\end{minipage}

\begin{remark}\label{rem:binary-coordinate-plane}{\rm
In Theorem~\ref{thm:binary-cz-kim}, the first coordinate plane
\(E_2=\mathbb F_2^2\) has Gram matrix \(I_2\). Its diagonal line
\(\mathbb F_2(1,1)\) is isotropic, but the form on \(E_2\) is
non-alternating; in particular, \(E_2\) is not an alternating hyperbolic
plane.

The arithmetic code is \(C=\iota_E(\operatorname{im}\Phi)\), with
\(\Phi\) and \(\iota_E\) defined above. The comparison with Kim's
matrix takes place in these Euclidean coordinates.

The coefficient group \(H^1_{\mathrm{et}}(U,\mu_2)\) is annihilated by
\(2\). If \(q\) is odd, every additive homomorphism
\(f:H^1_{\mathrm{et}}(U,\mu_2)\to\mathbb F_q^{2n}\) is therefore zero:
\[
2f(a)=f(2a)=0,\qquad 2\in\mathbb F_q^\times.
\]
Thus a nonzero odd-characteristic code cannot be obtained from this
restriction image by an additive coordinate identification. In
Theorem~\ref{thm:qary-building-up}, the correction terms instead lie on
\(\mathbb F_q(1,c)\), whose isotropy follows from \(c^2=-1\).
}
\end{remark}

\subsection{Normalized rank-one forms over odd fields}
\label{sec:odd-rank-one}

We return to \(K=\F_q\), \(q\equiv1\pmod4\), and fix \(c\in K^\times\)
with \(c^2=-1\).

To recover the following rank-one form, set \(k=n-1\), \(r=1\),
\(D=(1)\), and write the terminal row as \(\ell_i=(u_i,v_i)\).
The additional condition \(\ell_i\cdot\ell_i=-1\) forces the universal
diagonal coefficient \(c(1+\ell_i\cdot\ell_i)/2\) to vanish.  This
normalization is not a consequence of \(r=1\).

\begin{cor}[Normalized rank-one split boxed form over \texorpdfstring{$\F_q$}{Fq}]
\label{thm:split-boxed-form}
Fix \(c\in K^\times\) with \(c^2=-1\). For \(1\le i\le n-1\), let \(\ell_i=(u_i,v_i)\in K^2\) and \(a_i\in K\), and for distinct \(1\le i,j\le n-1\), let \(b_{ij}\in K\). Let \(G\) be the \(n\times 2n\) generator matrix whose \(n\times n\) block form \(\widetilde G=(B_{ij})\) of \(1\times 2\) blocks is determined by
\[
B_{ii}=(0,1)\quad (1\le i\le n-1),\qquad B_{nn}=(1,c),
\]
\[
B_{in}=\ell_i \quad (1\le i\le n-1),\qquad
B_{ni}=a_i(1,c)\quad (1\le i\le n-1),
\]
and
\[
B_{ij}=b_{ij}(1,c)\qquad (1\le i\neq j\le n-1).
\]
Thus \(\widetilde G\) has the split boxed shape
\[
\begingroup
\setlength{\arraycolsep}{9pt}
\renewcommand{\arraystretch}{1.72}
\widetilde G=
\left[
\begin{array}{c|cccc|c}
01 & b_{12}(1,c) & b_{13}(1,c) & \cdots & b_{1,n-1}(1,c) & \ell_1 \\
\hline
b_{21}(1,c) & 01 & b_{23}(1,c) & \cdots & b_{2,n-1}(1,c) & \ell_2 \\
b_{31}(1,c) & b_{32}(1,c) & 01 & \cdots & b_{3,n-1}(1,c) & \ell_3 \\
\vdots & \vdots & \vdots & \ddots & \vdots & \vdots \\
b_{n-1,1}(1,c) & b_{n-1,2}(1,c) & b_{n-1,3}(1,c) & \cdots & 01 & \ell_{n-1} \\
\hline
a_1(1,c) & a_2(1,c) & a_3(1,c) & \cdots & a_{n-1}(1,c) & (1,c)
\end{array}
\right],
\endgroup
\]

where each entry is a two-coordinate block.  Reading these blocks as scalar
pairs gives the ordinary \(n\times2n\) matrix \(G\).  Assume that
\begin{equation}
\label{eq:split-boxed-norm}
u_i^2+v_i^2=-1
\qquad (1\le i\le n-1),
\end{equation}
\begin{equation}
\label{eq:split-boxed-last-row}
u_i+c v_i=-c a_i
\qquad (1\le i\le n-1),
\end{equation}
and
\begin{equation}
\label{eq:split-boxed-offdiag}
c\,(b_{ij}+b_{ji})+\ell_i\cdot \ell_j=0
\qquad (1\le i<j\le n-1),
\end{equation}
where \(\ell_i\cdot \ell_j=u_i u_j+v_i v_j\) is the Euclidean dot product on the distinguished \(2\)-plane. Then the rows of \(G\) are pairwise orthogonal and linearly independent. Therefore \(G\) generates a self-dual code of length \(2n\).
\end{cor}

\begin{proof}
Take \(r=1\), \(D=(1)\), terminal row \(\ell_i=(u_i,v_i)\), and
off-diagonal coefficient \(b_{ij}\).  The lower-left coefficient forced by
Theorem~\ref{thm:universal-rank-boxed} is
\(-\partial_c\ell_i=cu_i-v_i\).  It equals the displayed coefficient
\(a_i\).  Since \(c^2=-1\), the equality
\(a_i=cu_i-v_i\) is equivalent to
\(u_i+cv_i=-ca_i\), which is
\eqref{eq:split-boxed-last-row}.  The diagonal entry of
the universal matrix is (01) precisely when its forced coefficient
\(c(1+\ell_i\cdot\ell_i)/2\) is zero, which is exactly
\eqref{eq:split-boxed-norm}.  For \(i\ne j\), the sole universal relation
is already \eqref{eq:split-boxed-offdiag}.  Thus the three hypotheses are
exactly the lower-left, diagonal, and off-diagonal conditions of the
rank-one universal form, and \(\det(1)=1\).
Theorem~\ref{thm:universal-rank-boxed} gives all three conclusions.

For clarity, the two non-final mixed contributions for rows \(i<j\) are
\[
(0,1)\cdot b_{ji}(1,c)+b_{ij}(1,c)\cdot(0,1)=cb_{ji}+cb_{ij};
\]
all other non-final off-diagonal contributions vanish since \((1,c)\cdot(1,c)=0\).
Linear independence also follows directly without any Gram hypothesis.
Writing \(g_i\) for the rows of \(G\), suppose that
\(\sum_i\lambda_i g_i=0\) and apply the linear functional
\(\beta(s,t)=t-cs\) to each non-final block column \(j\).
Only the diagonal block \((0,1)\) survives, so \(\lambda_j=0\).
The first coordinate of the final block then gives \(\lambda_n=0\).
\end{proof}

\noindent
In Lean, \(m=n-1\), the row and block index type is \texttt{Option (Fin m)}, and the split block-row space is canonically \(K^{2(m+1)}=K^{2n}\). The following formulation states pairwise orthogonality, linear independence, and self-duality under the three displayed coefficient conditions.

\noindent\begin{minipage}{\linewidth}
\begin{lstlisting}[caption={Lean for Corollary~\ref{thm:split-boxed-form}.}]
theorem paper_split_boxed_form_exact
    (c : K) (ell : Fin m -> SplitBlock K)
    (a : Fin m -> K) (b : Fin m -> Fin m -> K)
    (hc : c ^ 2 = (-1 : K))
    (hnorm : ∀ i, dot (ell i) (ell i) = (-1 : K))
    (hlast : ∀ i, ell i 0 + c * ell i 1 = -c * a i)
    (hoff : ∀ i j, i < j ->
      c * (b i j + b j i) + dot (ell i) (ell j) = 0) :
    SplitBoxedPairwiseOrthogonal (splitBoxedRows c ell a b) ∧
      LinearIndependent K (splitBoxedRows c ell a b) ∧
      splitBoxedRowSpace (splitBoxedRows c ell a b) =
        (splitBlockRowBilin (K := K) (m := m)).orthogonal
          (splitBoxedRowSpace (splitBoxedRows c ell a b))
\end{lstlisting}
\end{minipage}

\medskip

Conversely, self-duality forces the coefficient relations of
Corollary~\ref{thm:split-boxed-form} whenever a generator matrix satisfies
the following block conditions.

\begin{theorem}[Conditional boxed normalization from a distinguished isotropic-line row]
\label{thm:conditional-split-boxed-normalization}
Let \(n\geq1\), let \(C\subset K^{2n}\) be a self-dual code, and fix \(c\in K^\times\) with \(c^2=-1\). Assume that, after coordinate permutations preserving the \(2\)-block decomposition (hence Euclidean isometries) and elementary row operations, \(C\) admits a generator matrix of the following form. Work in these coordinates, and continue to denote the code by \(C\). Its last row \(g_n\) is
\[
\bigl(a_1(1,c),\,a_2(1,c),\,\dots,\,a_{n-1}(1,c),\,(1,c)\bigr).
\]
The remaining rows \(g_1,\dots,g_{n-1}\) satisfy the following block conditions:
\begin{enumerate}
\item the \(i\)-th block of \(g_i\) is \(01\);
\item the last block of \(g_i\) is \(\ell_i=(u_i,v_i)\);
\item every other non-final block of \(g_i\) lies on the same \(c\)-isotropic line, so for each \(j\neq i,n\) one has
\[
B_{ij}=b_{ij}(1,c)
\]
for some scalar \(b_{ij}\in K\).
\end{enumerate}
Then the coefficients satisfy
\[
u_i^2+v_i^2=-1,
\qquad
u_i+c v_i=-c a_i,
\qquad
c(b_{ij}+b_{ji})+\ell_i\cdot \ell_j=0,
\]
where the first two relations hold for \(1\leq i\leq n-1\) and the third
for \(1\leq i<j\leq n-1\). Thus the generator matrix is the matrix of
Corollary~\ref{thm:split-boxed-form}, with these coefficient relations.
\end{theorem}

\begin{proof}
Write the adapted generator matrix in block form as
\[
\widetilde G=(B_{ij})_{1\le i,j\le n},
\]
with \(B_{ii}=01\) for \(1\le i\le n-1\), \(B_{in}=\ell_i\), \(B_{ni}=a_i(1,c)\), and \(B_{ij}=b_{ij}(1,c)\) for \(i\neq j<n\). Since \(C\) is self-dual, its generator rows are pairwise orthogonal.

Applying self-orthogonality to a non-final row \(g_i\) gives
\[
0=\langle g_i,g_i\rangle
=
1+(u_i^2+v_i^2),
\]
because the diagonal block contributes \(1\), the final block contributes \(u_i^2+v_i^2\), and every other block is isotropic. Hence
\[
u_i^2+v_i^2=-1.
\]

Next, orthogonality of \(g_i\) with the distinguished last row gives
\[
0=\langle g_i,g_n\rangle
=
(0,1)\cdot a_i(1,c)+\ell_i\cdot(1,c)
=
c a_i+(u_i+c v_i),
\]
so
\[
u_i+c v_i=-c a_i.
\]

Finally, for \(1\le i<j\le n-1\), orthogonality of \(g_i\) and \(g_j\) gives
\[
\begin{aligned}
0=\langle g_i,g_j\rangle
&=(0,1)\cdot b_{ji}(1,c)
+b_{ij}(1,c)\cdot(0,1)+\ell_i\cdot\ell_j\\
&=c(b_{ij}+b_{ji})+\ell_i\cdot\ell_j,
\end{aligned}
\]
where all other block products vanish because \((1,c)\cdot(1,c)=0\).
These are exactly the boxed relations in Corollary~\ref{thm:split-boxed-form}.
\end{proof}

\noindent
In the following formulation, \(m=n-1\), and \texttt{none} denotes the
last row or block. The bilinear form is the Euclidean dot product in block
notation. The hypotheses specify the generator row space and every block
of the matrix; the conclusions give that matrix and its coefficient relations.

\noindent\begin{minipage}{\linewidth}
\begin{lstlisting}[caption={Lean formulation of Theorem~\ref{thm:conditional-split-boxed-normalization}.}]
theorem paper_conditional_boxed_normalization_exact
    {m : ℕ} (c : K) (hc : c ^ 2 = (-1 : K))
    (ell : Fin m → SplitBlock K) (a : Fin m → K)
    (b : Fin m → Fin m → K)
    (G : Option (Fin m) → SplitBlockRow K m)
    (C : Submodule K (SplitBlockRow K m))
    (hC : C = (splitBlockRowBilin (K := K) (m := m)).orthogonal C)
    (hgen : splitBoxedRowSpace G = C)
    (hfinal : G none none = isotropicLineBlock c 1)
    (hfinalLeft : ∀ j, G none (some j) = isotropicLineBlock c (a j))
    (hlast : ∀ i, G (some i) none = ell i)
    (hdiag : ∀ i, G (some i) (some i) = splitDiagonalBlock)
    (hother : ∀ i j, i ≠ j →
      G (some i) (some j) = isotropicLineBlock c (b i j)) :
    G = splitBoxedRows c ell a b ∧
      (∀ i, dot (ell i) (ell i) = (-1 : K)) ∧
      (∀ i, ell i 0 + c * ell i 1 = -c * a i) ∧
      (∀ i j, i < j → c * (b i j + b j i) + dot (ell i) (ell j) = 0) ∧
      C = splitBoxedRowSpace (splitBoxedRows c ell a b)
\end{lstlisting}
\end{minipage}


\begin{remark}
\label{rem:universal-and-rank-one-scope}
\rm
Let \(K=\mathbb F_q\), where \(q\equiv1\pmod4\), choose \(c\in K\)
with \(c^2=-1\), and fix a pairing of the \(2n\) coordinates.
Put \(U_c=\bigoplus_{j=1}^n K(1,c)\). For every
\(C=C^\perp\subseteq K^{2n}\), Theorem~\ref{thm:universal-rank-boxed}
gives a generator \(G(c;b,\ell,D)\), after a permutation of the
two-coordinate blocks, with
\[
r=\dim(C\cap U_c),\qquad k=n-r,\qquad \det D\ne0.
\]
When \(r=1\), scaling the terminal row normalizes \(D\) to \((1)\).
The further conditions \(\ell_i\cdot\ell_i=-1\) are additional assumptions,
not consequences of \(r=1\). Under these conditions, the universal
relations reduce directly to
\eqref{eq:split-boxed-norm}--\eqref{eq:split-boxed-offdiag}, giving
Corollary~\ref{thm:split-boxed-form}. For a generator satisfying all the
block conditions of Theorem~\ref{thm:conditional-split-boxed-normalization},
self-duality implies precisely these three coefficient relations.
\end{remark}

\section{Applications}\label{sec:applications}

% Generated by build_application_catalog.py; do not edit.
\newcommand{\ApplicationCatalogueRows}{%
\href{https://github.com/LeGenAI/intersection-coding-theory-cohomology/blob/afm-revision-2026-08-27/Formalization/Verification/Examples/certificates/gf5-06.json}{\texttt{GF5-06}} & 5 & $[6,3,4];\,60$ & $(4;\,60)$ & Leon--Pless--Sloane database \\
\href{https://github.com/LeGenAI/intersection-coding-theory-cohomology/blob/afm-revision-2026-08-27/Formalization/Verification/Examples/certificates/gf5-08.json}{\texttt{GF5-08}} & 5 & $[8,4,4];\,48$ & $(4;\,48)$ & Leon--Pless--Sloane database \\
\href{https://github.com/LeGenAI/intersection-coding-theory-cohomology/blob/afm-revision-2026-08-27/Formalization/Verification/Examples/certificates/gf13-08.json}{\texttt{GF13-08}} & 13 & $[8,4,5];\,672$ & $(5;\,672)$ & Betsumiya et al. \\
\href{https://github.com/LeGenAI/intersection-coding-theory-cohomology/blob/afm-revision-2026-08-27/Formalization/Verification/Examples/certificates/gf13-10.json}{\texttt{GF13-10}} & 13 & $[10,5,6];\,2,520$ & $(6;\,2,520)$ & Betsumiya et al. \\
\href{https://github.com/LeGenAI/intersection-coding-theory-cohomology/blob/afm-revision-2026-08-27/Formalization/Verification/Examples/certificates/gf13-12a.json}{\texttt{GF13-12A}} & 13 & $[12,6,6];\,576$ & $(6;\,\{528,576,696,792\})$ & Betsumiya et al. \\
\href{https://github.com/LeGenAI/intersection-coding-theory-cohomology/blob/afm-revision-2026-08-27/Formalization/Verification/Examples/certificates/gf13-12p.json}{\texttt{GF13-12P}} & 13 & $[12,6,6];\,960$ & $(6;\,\{528,576,696,792\})$ & Betsumiya et al. \\
\href{https://github.com/LeGenAI/intersection-coding-theory-cohomology/blob/afm-revision-2026-08-27/Formalization/Verification/Examples/certificates/gf13-14.json}{\texttt{GF13-14}} & 13 & $[14,7,8];\,36,036$ & $(8;\,36,036)$ & Betsumiya et al. \\
}

% Generated by check_large_applications.py; do not edit.
\newcommand{\AppThirteenFourteenUniversal}{%
\left[\begin{array}{cccccc|c}
(0,1) & (12,8) & (11,3) & (5,12) & (6,4) & (9,6) & (3,9)\\
(1,5) & (0,1) & (12,8) & (11,3) & (4,7) & (1,5) & (9,10)\\
(10,11) & (1,5) & (0,1) & (3,2) & (11,3) & (0,0) & (3,9)\\
(8,1) & (7,9) & (10,11) & (0,1) & (9,6) & (9,6) & (4,3)\\
(3,2) & (6,4) & (11,3) & (7,9) & (0,1) & (9,6) & (0,8)\\
(12,8) & (12,8) & (8,1) & (4,7) & (0,0) & (0,1) & (3,9)\\
\hline
(6,4) & (9,6) & (6,4) & (4,7) & (5,12) & (6,4) & (1,5)\\
\end{array}\right]
}
\newcommand{\AppThirteenFourteenP}{%
\begin{pmatrix}
0 & 12 & 11 & 5 & 6 & 9\\
1 & 0 & 12 & 11 & 4 & 1\\
10 & 1 & 0 & 3 & 11 & 0\\
8 & 7 & 10 & 0 & 9 & 9\\
3 & 6 & 11 & 7 & 0 & 9\\
12 & 12 & 8 & 4 & 0 & 0\\
\end{pmatrix}
}
\newcommand{\AppThirteenFourteenH}{%
\begin{pmatrix}
3\\
9\\
3\\
4\\
0\\
3\\
\end{pmatrix}
}
\newcommand{\AppThirteenFourteenQ}{%
\begin{pmatrix}
7\\
4\\
7\\
9\\
8\\
7\\
\end{pmatrix}
}
\newcommand{\AppThirteenFourteenA}{%
\begin{pmatrix}
6 & 9 & 6 & 4 & 5 & 6\\
\end{pmatrix}
}
\newcommand{\AppThirteenFourteenReadout}{%
\left[\begin{array}{cccccc|c}
1 & 0 & 0 & 0 & 0 & 0 & 7\\
0 & 1 & 0 & 0 & 0 & 0 & 4\\
0 & 0 & 1 & 0 & 0 & 0 & 7\\
0 & 0 & 0 & 1 & 0 & 0 & 9\\
0 & 0 & 0 & 0 & 1 & 0 & 8\\
0 & 0 & 0 & 0 & 0 & 1 & 7\\
\hline
0 & 0 & 0 & 0 & 0 & 0 & 0\\
\end{array}\right]
}
\newcommand{\AppThirteenFourteenGamma}{%
(1,10,8,3,12,6)
}
\newcommand{\AppThirteenFourteenSurvivor}{%
[1:10:8:3:12:6]
}
\newcommand{\AppThirteenFourteenCoordinateOrder}{%
(13,9,1,10,14,6,8,2,7,4,12,3,5,11)
}
\newcommand{\AppThirteenFourteenTransformDeterminant}{%
7
}
\newcommand{\AppThirteenFourteenDistanceRow}{%
\ref{prop:universal-mds-137} & 13 & $[14,7]$ & 62,748,516 & 8 & 36,036 \\
}
\newcommand{\AppThirteenFourteenParentDistribution}{%
(1,0,0,0,0,0,960,3744,50040,226320,853920,1844064,1847760)
}
\newcommand{\AppThirteenFourteenLowLines}{%
392
}

% Generated by check_golay_lineage.py; do not edit.
\newcommand{\GolayLineageRows}{%
0 & $C_{16}$ & $[16,8,4]$ & $\mathbf{12}$ & base parent \\
1 & $C_{18}$ & $[18,9,4]$ & $\mathbf{9}$ & $\mathcal B_1(G_{16},x_{16})$ \\
2 & $C_{20}$ & $[20,10,4]$ & $\mathbf{5}$ & $\mathcal B_1(G_{18},x_{18})$ \\
3 & $C_{22}$ & $[22,11,6]$ & $\mathbf{77}$ & $\mathcal B_1(G_{20},x_{20})$ \\
4 & $\mathcal G_{24}$ & $[24,12,8]$ & $\mathbf{759}$ & $\mathcal B_1(G_{22},x_{22})$ \\
}
\newcommand{\GolayCatalogueRows}{%
\href{https://github.com/LeGenAI/intersection-coding-theory-cohomology/blob/afm-revision-2026-08-27/Formalization/Verification/Examples/certificates/binary-golay-lineage.json}{$C_{16}^{(2)}$} & 2 & $[16,8,4];\,\mathbf{12}$ & $(4;\,12)$ & Harada--Munemasa database \\
\href{https://github.com/LeGenAI/intersection-coding-theory-cohomology/blob/afm-revision-2026-08-27/Formalization/Verification/Examples/certificates/binary-golay-lineage.json}{$C_{18}^{(2)}$} & 2 & $[18,9,4];\,\mathbf{9}$ & $(4;\,9)$ & Harada--Munemasa database \\
\href{https://github.com/LeGenAI/intersection-coding-theory-cohomology/blob/afm-revision-2026-08-27/Formalization/Verification/Examples/certificates/binary-golay-lineage.json}{$C_{20}^{(2)}$} & 2 & $[20,10,4];\,\mathbf{5}$ & $(4;\,5)$ & Harada--Munemasa database \\
\href{https://github.com/LeGenAI/intersection-coding-theory-cohomology/blob/afm-revision-2026-08-27/Formalization/Verification/Examples/certificates/binary-golay-lineage.json}{$C_{22}^{(2)}$} & 2 & $[22,11,6];\,\mathbf{77}$ & $(6;\,77)$ & Harada--Munemasa database \\
\href{https://github.com/LeGenAI/intersection-coding-theory-cohomology/blob/afm-revision-2026-08-27/Formalization/Verification/Examples/certificates/binary-golay-lineage.json}{$\mathcal G_{24}$} & 2 & $[24,12,8];\,\mathbf{759}$ & $(8;\,759)$ & Harada--Munemasa database \\
}
\newcommand{\BinaryLargestRepeatedMatrix}{\left(\begin{array}{r|cc!{\vrule width 1.2pt}cc|cc|cc|cc|cc|cc|cc|cc|cc|cc!{\vrule width 1.2pt}cc}
\scriptstyle\text{rows} & \multicolumn{2}{c!{\vrule width 1.2pt}}{\scriptstyle\text{new pair}} & \multicolumn{20}{c!{\vrule width 1.2pt}}{\scriptstyle \operatorname{diag}:\ 01;\quad i\ne j:\ b_{ij}(11)} & \multicolumn{2}{c}{\scriptstyle D=(1):\ 10;\ 11} \\ \noalign{\hrule height 0.7pt}
\scriptstyle b_{0j},10 & \cellcolor{black!7}\textcolor{blue!65!black}{0} & \cellcolor{black!7}\textcolor{blue!65!black}{1} & \textcolor{blue!65!black}{0} & \textcolor{blue!65!black}{0} & \textcolor{blue!65!black}{1} & \textcolor{blue!65!black}{1} & \textcolor{blue!65!black}{0} & \textcolor{blue!65!black}{0} & \textcolor{blue!65!black}{1} & \textcolor{blue!65!black}{1} & \textcolor{blue!65!black}{1} & \textcolor{blue!65!black}{1} & \textcolor{blue!65!black}{1} & \textcolor{blue!65!black}{1} & \textcolor{blue!65!black}{0} & \textcolor{blue!65!black}{0} & \textcolor{blue!65!black}{0} & \textcolor{blue!65!black}{0} & \textcolor{blue!65!black}{0} & \textcolor{blue!65!black}{0} & \textcolor{blue!65!black}{1} & \textcolor{blue!65!black}{1} & \textcolor{blue!65!black}{1} & \textcolor{blue!65!black}{0} \\ \noalign{\hrule height 1.2pt}
\scriptstyle b_{ij},10 & \textcolor{orange!80!black}{1} & \textcolor{orange!80!black}{1} & \cellcolor{black!7}\textcolor{teal!60!black}{0} & \cellcolor{black!7}\textcolor{teal!60!black}{1} & \textcolor{teal!60!black}{0} & \textcolor{teal!60!black}{0} & \textcolor{teal!60!black}{1} & \textcolor{teal!60!black}{1} & \textcolor{teal!60!black}{0} & \textcolor{teal!60!black}{0} & \textcolor{teal!60!black}{1} & \textcolor{teal!60!black}{1} & \textcolor{teal!60!black}{1} & \textcolor{teal!60!black}{1} & \textcolor{teal!60!black}{1} & \textcolor{teal!60!black}{1} & \textcolor{teal!60!black}{0} & \textcolor{teal!60!black}{0} & \textcolor{teal!60!black}{0} & \textcolor{teal!60!black}{0} & \textcolor{teal!60!black}{0} & \textcolor{teal!60!black}{0} & \textcolor{teal!60!black}{1} & \textcolor{teal!60!black}{0} \\
\scriptstyle b_{ij},10 & \textcolor{orange!80!black}{0} & \textcolor{orange!80!black}{0} & \textcolor{teal!60!black}{1} & \textcolor{teal!60!black}{1} & \cellcolor{black!7}\textcolor{teal!60!black}{0} & \cellcolor{black!7}\textcolor{teal!60!black}{1} & \textcolor{teal!60!black}{0} & \textcolor{teal!60!black}{0} & \textcolor{teal!60!black}{0} & \textcolor{teal!60!black}{0} & \textcolor{teal!60!black}{1} & \textcolor{teal!60!black}{1} & \textcolor{teal!60!black}{0} & \textcolor{teal!60!black}{0} & \textcolor{teal!60!black}{1} & \textcolor{teal!60!black}{1} & \textcolor{teal!60!black}{1} & \textcolor{teal!60!black}{1} & \textcolor{teal!60!black}{0} & \textcolor{teal!60!black}{0} & \textcolor{teal!60!black}{1} & \textcolor{teal!60!black}{1} & \textcolor{teal!60!black}{1} & \textcolor{teal!60!black}{0} \\
\scriptstyle b_{ij},10 & \textcolor{orange!80!black}{1} & \textcolor{orange!80!black}{1} & \textcolor{teal!60!black}{0} & \textcolor{teal!60!black}{0} & \textcolor{teal!60!black}{1} & \textcolor{teal!60!black}{1} & \cellcolor{black!7}\textcolor{teal!60!black}{0} & \cellcolor{black!7}\textcolor{teal!60!black}{1} & \textcolor{teal!60!black}{0} & \textcolor{teal!60!black}{0} & \textcolor{teal!60!black}{0} & \textcolor{teal!60!black}{0} & \textcolor{teal!60!black}{1} & \textcolor{teal!60!black}{1} & \textcolor{teal!60!black}{0} & \textcolor{teal!60!black}{0} & \textcolor{teal!60!black}{1} & \textcolor{teal!60!black}{1} & \textcolor{teal!60!black}{1} & \textcolor{teal!60!black}{1} & \textcolor{teal!60!black}{0} & \textcolor{teal!60!black}{0} & \textcolor{teal!60!black}{1} & \textcolor{teal!60!black}{0} \\
\scriptstyle b_{ij},10 & \textcolor{orange!80!black}{0} & \textcolor{orange!80!black}{0} & \textcolor{teal!60!black}{1} & \textcolor{teal!60!black}{1} & \textcolor{teal!60!black}{1} & \textcolor{teal!60!black}{1} & \textcolor{teal!60!black}{1} & \textcolor{teal!60!black}{1} & \cellcolor{black!7}\textcolor{teal!60!black}{0} & \cellcolor{black!7}\textcolor{teal!60!black}{1} & \textcolor{teal!60!black}{1} & \textcolor{teal!60!black}{1} & \textcolor{teal!60!black}{1} & \textcolor{teal!60!black}{1} & \textcolor{teal!60!black}{1} & \textcolor{teal!60!black}{1} & \textcolor{teal!60!black}{0} & \textcolor{teal!60!black}{0} & \textcolor{teal!60!black}{1} & \textcolor{teal!60!black}{1} & \textcolor{teal!60!black}{0} & \textcolor{teal!60!black}{0} & \textcolor{teal!60!black}{1} & \textcolor{teal!60!black}{0} \\
\scriptstyle b_{ij},10 & \textcolor{orange!80!black}{0} & \textcolor{orange!80!black}{0} & \textcolor{teal!60!black}{0} & \textcolor{teal!60!black}{0} & \textcolor{teal!60!black}{0} & \textcolor{teal!60!black}{0} & \textcolor{teal!60!black}{1} & \textcolor{teal!60!black}{1} & \textcolor{teal!60!black}{0} & \textcolor{teal!60!black}{0} & \cellcolor{black!7}\textcolor{teal!60!black}{0} & \cellcolor{black!7}\textcolor{teal!60!black}{1} & \textcolor{teal!60!black}{0} & \textcolor{teal!60!black}{0} & \textcolor{teal!60!black}{1} & \textcolor{teal!60!black}{1} & \textcolor{teal!60!black}{1} & \textcolor{teal!60!black}{1} & \textcolor{teal!60!black}{0} & \textcolor{teal!60!black}{0} & \textcolor{teal!60!black}{0} & \textcolor{teal!60!black}{0} & \textcolor{teal!60!black}{1} & \textcolor{teal!60!black}{0} \\
\scriptstyle b_{ij},10 & \textcolor{orange!80!black}{0} & \textcolor{orange!80!black}{0} & \textcolor{teal!60!black}{0} & \textcolor{teal!60!black}{0} & \textcolor{teal!60!black}{1} & \textcolor{teal!60!black}{1} & \textcolor{teal!60!black}{0} & \textcolor{teal!60!black}{0} & \textcolor{teal!60!black}{0} & \textcolor{teal!60!black}{0} & \textcolor{teal!60!black}{1} & \textcolor{teal!60!black}{1} & \cellcolor{black!7}\textcolor{teal!60!black}{0} & \cellcolor{black!7}\textcolor{teal!60!black}{1} & \textcolor{teal!60!black}{0} & \textcolor{teal!60!black}{0} & \textcolor{teal!60!black}{1} & \textcolor{teal!60!black}{1} & \textcolor{teal!60!black}{1} & \textcolor{teal!60!black}{1} & \textcolor{teal!60!black}{1} & \textcolor{teal!60!black}{1} & \textcolor{teal!60!black}{1} & \textcolor{teal!60!black}{0} \\
\scriptstyle b_{ij},10 & \textcolor{orange!80!black}{1} & \textcolor{orange!80!black}{1} & \textcolor{teal!60!black}{0} & \textcolor{teal!60!black}{0} & \textcolor{teal!60!black}{0} & \textcolor{teal!60!black}{0} & \textcolor{teal!60!black}{1} & \textcolor{teal!60!black}{1} & \textcolor{teal!60!black}{0} & \textcolor{teal!60!black}{0} & \textcolor{teal!60!black}{0} & \textcolor{teal!60!black}{0} & \textcolor{teal!60!black}{1} & \textcolor{teal!60!black}{1} & \cellcolor{black!7}\textcolor{teal!60!black}{0} & \cellcolor{black!7}\textcolor{teal!60!black}{1} & \textcolor{teal!60!black}{0} & \textcolor{teal!60!black}{0} & \textcolor{teal!60!black}{1} & \textcolor{teal!60!black}{1} & \textcolor{teal!60!black}{1} & \textcolor{teal!60!black}{1} & \textcolor{teal!60!black}{1} & \textcolor{teal!60!black}{0} \\
\scriptstyle b_{ij},10 & \textcolor{orange!80!black}{1} & \textcolor{orange!80!black}{1} & \textcolor{teal!60!black}{1} & \textcolor{teal!60!black}{1} & \textcolor{teal!60!black}{0} & \textcolor{teal!60!black}{0} & \textcolor{teal!60!black}{0} & \textcolor{teal!60!black}{0} & \textcolor{teal!60!black}{1} & \textcolor{teal!60!black}{1} & \textcolor{teal!60!black}{0} & \textcolor{teal!60!black}{0} & \textcolor{teal!60!black}{0} & \textcolor{teal!60!black}{0} & \textcolor{teal!60!black}{1} & \textcolor{teal!60!black}{1} & \cellcolor{black!7}\textcolor{teal!60!black}{0} & \cellcolor{black!7}\textcolor{teal!60!black}{1} & \textcolor{teal!60!black}{0} & \textcolor{teal!60!black}{0} & \textcolor{teal!60!black}{1} & \textcolor{teal!60!black}{1} & \textcolor{teal!60!black}{1} & \textcolor{teal!60!black}{0} \\
\scriptstyle b_{ij},10 & \textcolor{orange!80!black}{1} & \textcolor{orange!80!black}{1} & \textcolor{teal!60!black}{1} & \textcolor{teal!60!black}{1} & \textcolor{teal!60!black}{1} & \textcolor{teal!60!black}{1} & \textcolor{teal!60!black}{0} & \textcolor{teal!60!black}{0} & \textcolor{teal!60!black}{0} & \textcolor{teal!60!black}{0} & \textcolor{teal!60!black}{1} & \textcolor{teal!60!black}{1} & \textcolor{teal!60!black}{0} & \textcolor{teal!60!black}{0} & \textcolor{teal!60!black}{0} & \textcolor{teal!60!black}{0} & \textcolor{teal!60!black}{1} & \textcolor{teal!60!black}{1} & \cellcolor{black!7}\textcolor{teal!60!black}{0} & \cellcolor{black!7}\textcolor{teal!60!black}{1} & \textcolor{teal!60!black}{0} & \textcolor{teal!60!black}{0} & \textcolor{teal!60!black}{1} & \textcolor{teal!60!black}{0} \\
\scriptstyle b_{ij},10 & \textcolor{orange!80!black}{0} & \textcolor{orange!80!black}{0} & \textcolor{teal!60!black}{1} & \textcolor{teal!60!black}{1} & \textcolor{teal!60!black}{0} & \textcolor{teal!60!black}{0} & \textcolor{teal!60!black}{1} & \textcolor{teal!60!black}{1} & \textcolor{teal!60!black}{1} & \textcolor{teal!60!black}{1} & \textcolor{teal!60!black}{1} & \textcolor{teal!60!black}{1} & \textcolor{teal!60!black}{0} & \textcolor{teal!60!black}{0} & \textcolor{teal!60!black}{0} & \textcolor{teal!60!black}{0} & \textcolor{teal!60!black}{0} & \textcolor{teal!60!black}{0} & \textcolor{teal!60!black}{1} & \textcolor{teal!60!black}{1} & \cellcolor{black!7}\textcolor{teal!60!black}{0} & \cellcolor{black!7}\textcolor{teal!60!black}{1} & \textcolor{teal!60!black}{1} & \textcolor{teal!60!black}{0} \\ \noalign{\hrule height 1.2pt}
\scriptstyle 11,11 & \textcolor{orange!80!black}{1} & \textcolor{orange!80!black}{1} & \textcolor{violet!75!black}{1} & \textcolor{violet!75!black}{1} & \textcolor{violet!75!black}{1} & \textcolor{violet!75!black}{1} & \textcolor{violet!75!black}{1} & \textcolor{violet!75!black}{1} & \textcolor{violet!75!black}{1} & \textcolor{violet!75!black}{1} & \textcolor{violet!75!black}{1} & \textcolor{violet!75!black}{1} & \textcolor{violet!75!black}{1} & \textcolor{violet!75!black}{1} & \textcolor{violet!75!black}{1} & \textcolor{violet!75!black}{1} & \textcolor{violet!75!black}{1} & \textcolor{violet!75!black}{1} & \textcolor{violet!75!black}{1} & \textcolor{violet!75!black}{1} & \textcolor{violet!75!black}{1} & \textcolor{violet!75!black}{1} & \cellcolor{black!7}\textcolor{violet!75!black}{1} & \cellcolor{black!7}\textcolor{violet!75!black}{1} \\
\end{array}\right)}
\newcommand{\BinaryLargestParentParameters}{c=1,\; k=10,\; r=1,\; D=\begin{pmatrix}1\end{pmatrix}}
\newcommand{\GolaySixteenCorrection}{\texttt{0000001100110111}}
\newcommand{\GolayEighteenCorrection}{\texttt{000000110011011001}}
\newcommand{\GolayTwentyCorrection}{\texttt{00000000110100101011}}
\newcommand{\GolayTwentyTwoCorrection}{\texttt{0000000000101011100011}}
\newcommand{\GolayLineageCertificate}{\href{https://github.com/LeGenAI/intersection-coding-theory-cohomology/blob/afm-revision-2026-08-27/Formalization/Verification/Examples/certificates/binary-golay-lineage.json}{$\mathcal G_{24}$}}

% Generated by check_gf5_repeated_top.py; do not edit.
\newcommand{\GFFiveRepeatedLineageRows}{%
0 & $C_{22}^{(5)}$ & $[22,11,8]$ & 660 & base parent \\
1 & $C_{24}^{(5)}$ & $[24,12,9]$ & 1,056 & $\mathcal B_2(G_{22}^{(5)},x_{22}^{(5)})$ \\
}
\newcommand{\GFFiveRepeatedCatalogueRows}{%
\href{https://github.com/LeGenAI/intersection-coding-theory-cohomology/blob/afm-revision-2026-08-27/Formalization/Verification/Examples/certificates/gf5-repeated-top.json}{$C_{24}^{(5)}$} & 5 & $[24,12,9];\,1,056$ & $(9;\,1,056)$ & Georgiou--Koukouvinos \\
}
\newcommand{\GFFiveLargestRepeatedMatrix}{\left(\begin{array}{cc|cc|cc|cc|cc|cc|cc|cc|cc|cc|cc|cc}
1 & 0 & 1 & 0 & 0 & 0 & 0 & 0 & 0 & 0 & 0 & 0 & 0 & 0 & 1 & 1 & 1 & 1 & 1 & 0 & 0 & 3 & 0 & 2 \\
1 & 2 & 2 & 0 & 0 & 0 & 0 & 0 & 0 & 0 & 0 & 0 & 0 & 3 & 4 & 4 & 4 & 4 & 4 & 2 & 3 & 1 & 2 & 2 \\
4 & 3 & 0 & 2 & 0 & 0 & 0 & 0 & 0 & 0 & 0 & 0 & 0 & 4 & 1 & 3 & 3 & 3 & 1 & 3 & 4 & 4 & 2 & 4 \\
1 & 2 & 0 & 0 & 2 & 0 & 0 & 0 & 0 & 0 & 0 & 0 & 0 & 1 & 2 & 2 & 4 & 4 & 4 & 3 & 1 & 4 & 2 & 3 \\
1 & 2 & 0 & 0 & 0 & 2 & 0 & 0 & 0 & 0 & 0 & 0 & 0 & 1 & 2 & 4 & 2 & 4 & 3 & 2 & 3 & 4 & 4 & 1 \\
1 & 2 & 0 & 0 & 0 & 0 & 2 & 0 & 0 & 0 & 0 & 0 & 0 & 1 & 2 & 4 & 4 & 2 & 2 & 4 & 2 & 1 & 2 & 1 \\
1 & 2 & 0 & 0 & 0 & 0 & 0 & 2 & 0 & 0 & 0 & 0 & 0 & 1 & 4 & 2 & 3 & 4 & 2 & 4 & 3 & 1 & 4 & 3 \\
4 & 3 & 0 & 0 & 0 & 0 & 0 & 0 & 2 & 0 & 0 & 0 & 0 & 2 & 1 & 2 & 3 & 1 & 3 & 1 & 4 & 1 & 3 & 4 \\
1 & 2 & 0 & 0 & 0 & 0 & 0 & 0 & 0 & 2 & 0 & 0 & 0 & 3 & 2 & 4 & 2 & 3 & 4 & 4 & 1 & 1 & 4 & 3 \\
1 & 2 & 0 & 0 & 0 & 0 & 0 & 0 & 0 & 0 & 2 & 0 & 0 & 2 & 4 & 4 & 4 & 2 & 2 & 2 & 1 & 4 & 4 & 3 \\
4 & 3 & 0 & 0 & 0 & 0 & 0 & 0 & 0 & 0 & 0 & 2 & 0 & 2 & 2 & 3 & 1 & 3 & 1 & 1 & 2 & 1 & 1 & 4 \\
0 & 0 & 0 & 0 & 0 & 0 & 0 & 0 & 0 & 0 & 0 & 0 & 2 & 1 & 4 & 4 & 1 & 1 & 4 & 4 & 4 & 1 & 1 & 4
\end{array}\right)}
\newcommand{\GFFiveTwentyTwoCorrection}{(1,0,0,0,0,0,0,0,0,0,0,0,1,1,1,1,1,0,0,3,0,2)}
\newcommand{\GFFiveRepeatedCertificate}{\href{https://github.com/LeGenAI/intersection-coding-theory-cohomology/blob/afm-revision-2026-08-27/Formalization/Verification/Examples/certificates/gf5-repeated-top.json}{$C_{24}^{(5)}$}}

% Generated by check_gf13_repeated_lineage.py; do not edit.
\newcommand{\GFThirteenRepeatedLineageRows}{%
0 & $C_{18}^{(13)}$ & $[18,9,7]$ & 72 & base parent \\
1 & $C_{20}^{(13)}$ & $[20,10,8]$ & 120 & $\mathcal B_5(G_{18}^{(13)},x_{18}^{(13)})$ \\
2 & $C_{22}^{(13)}$ & $[22,11,10]$ & 23,628 & $\mathcal B_5(G_{20}^{(13)},x_{20}^{(13)})$ \\
}
\newcommand{\GFThirteenRepeatedCatalogueRows}{%
\href{https://github.com/LeGenAI/intersection-coding-theory-cohomology/blob/afm-revision-2026-08-27/Formalization/Verification/Examples/certificates/gf13-repeated-lineage.json}{$C_{18}^{(13)}$} & 13 & $[18,9,7];\,72$ & $(8;\,2,484)$ & Betsumiya et al. \\
\href{https://github.com/LeGenAI/intersection-coding-theory-cohomology/blob/afm-revision-2026-08-27/Formalization/Verification/Examples/certificates/gf13-repeated-lineage.json}{$C_{20}^{(13)}$} & 13 & $[20,10,8];\,120$ & $(10;\,111,024)$ & Betsumiya et al. \\
\href{https://github.com/LeGenAI/intersection-coding-theory-cohomology/blob/afm-revision-2026-08-27/Formalization/Verification/Examples/certificates/gf13-repeated-lineage.json}{$C_{22}^{(13)}$} & 13 & $[22,11,10];\,23,628$ & $(10;\,23,628)$ & Betsumiya et al. \\
}
\newcommand{\GFThirteenLargestRepeatedMatrix}{\left(\begin{array}{cc|cc|cc|cc|cc|cc|cc|cc|cc|cc|cc}
1 & 0 & 0 & 0 & 0 & 0 & 0 & 0 & 0 & 0 & 0 & 0 & 7 & 7 & 8 & 3 & 10 & 1 & 9 & 7 & 11 & 10 \\
11 & 3 & 1 & 0 & 0 & 0 & 0 & 0 & 0 & 0 & 0 & 0 & 9 & 6 & 4 & 2 & 6 & 8 & 9 & 8 & 11 & 4 \\
12 & 8 & 0 & 1 & 0 & 0 & 0 & 0 & 0 & 0 & 0 & 0 & 4 & 3 & 12 & 4 & 6 & 12 & 1 & 7 & 11 & 10 \\
4 & 7 & 0 & 0 & 1 & 0 & 0 & 0 & 0 & 0 & 0 & 0 & 9 & 0 & 3 & 6 & 9 & 7 & 10 & 12 & 6 & 10 \\
4 & 7 & 0 & 0 & 0 & 1 & 0 & 0 & 0 & 0 & 0 & 0 & 11 & 9 & 4 & 9 & 8 & 12 & 0 & 2 & 2 & 2 \\
2 & 10 & 0 & 0 & 0 & 0 & 1 & 0 & 0 & 0 & 0 & 0 & 2 & 10 & 10 & 4 & 1 & 9 & 0 & 4 & 9 & 4 \\
10 & 11 & 0 & 0 & 0 & 0 & 0 & 1 & 0 & 0 & 0 & 0 & 2 & 6 & 11 & 11 & 6 & 4 & 6 & 0 & 9 & 9 \\
4 & 7 & 0 & 0 & 0 & 0 & 0 & 0 & 1 & 0 & 0 & 0 & 2 & 12 & 7 & 8 & 12 & 1 & 4 & 9 & 0 & 4 \\
7 & 9 & 0 & 0 & 0 & 0 & 0 & 0 & 0 & 1 & 0 & 0 & 0 & 10 & 1 & 9 & 12 & 5 & 8 & 4 & 6 & 0 \\
8 & 1 & 0 & 0 & 0 & 0 & 0 & 0 & 0 & 0 & 1 & 0 & 12 & 7 & 12 & 8 & 11 & 2 & 5 & 1 & 4 & 9 \\
2 & 10 & 0 & 0 & 0 & 0 & 0 & 0 & 0 & 0 & 0 & 1 & 8 & 9 & 6 & 6 & 8 & 11 & 12 & 12 & 6 & 1
\end{array}\right)}
\newcommand{\GFThirteenEighteenCorrection}{(0,0,0,0,0,0,0,1,4,2,10,8,1,6,8,7,4,5)}
\newcommand{\GFThirteenTwentyCorrection}{(0,0,0,0,0,0,0,0,0,0,7,7,8,3,10,1,9,7,11,10)}
\newcommand{\GFThirteenRepeatedCertificate}{\href{https://github.com/LeGenAI/intersection-coding-theory-cohomology/blob/afm-revision-2026-08-27/Formalization/Verification/Examples/certificates/gf13-repeated-lineage.json}{$C_{22}^{(13)}$}}


The applications begin with an explicit realization of an optimal code in
the universal form and the correction search exposed by its exact parent.
The comparison table then places the resulting codes against published
benchmarks, and the final subsection displays the largest repeated
parent--child realizations.

\subsection{Universal realization and correction search over
\texorpdfstring{$\GF{13}$}{GF(13)}}

The first application starts from an optimal code and asks whether the
universal form exposes a useful parent.  The answer is exact: a known MDS
code has a rank-one realization, its two-coordinate reduction is self-dual,
and only one projective correction class over that parent can restore the MDS
distance.

\begin{prop}
\label{prop:universal-mds-137}
\rm
Let \(K=\GF{13}\), let \(c=5\), and let
\[
G_{14}=\left[I_7\ \middle|\
\operatorname{circ}(2,9,10,9,2,1,1)\right].
\]
This is the pure double-circulant self-dual MDS \([14,7,8]\) code recorded
in~\cite{BetsumiyaEtAl}. Put the columns of \(G_{14}\) in the order
\[
\sigma=\AppThirteenFourteenCoordinateOrder
\]
and group consecutive coordinates into seven two-coordinate blocks. A
row-equivalent universal generator is obtained from
\begin{equation}
P=\AppThirteenFourteenP,
\qquad
H=\AppThirteenFourteenH,
\qquad
Q=\AppThirteenFourteenQ.
\label{eq:gf13-14-coefficients}
\end{equation}
Let \(\delta_{ij}\) denote the Kronecker delta.  In the notation of
Theorem~\ref{thm:universal-rank-boxed},
\begin{equation}
\widehat G_{14}=G(5;P,H,Q,(1))
=\left[
\begin{array}{c|c}
\bigl((P_{ij},5P_{ij}+\delta_{ij})\bigr)_{1\le i,j\le6}
  & \bigl((H_i,5H_i+Q_i)\bigr)_{1\le i\le6}\\
\hline
\bigl(-Q_j(1,5)\bigr)_{1\le j\le6} & (1,5)
\end{array}
\right].
\label{eq:gf13-14-structured-generator}
\end{equation}
It has \(k=6\), \(r=1\), and \(D=(1)\).  Deleting its first block row and
first block column gives a self-dual \([12,6,6]\) parent.  For this parent,
the projective correction classes satisfying the necessary weight-six/seven
lifting condition form the singleton
\[
\Lambda_8=\{\AppThirteenFourteenSurvivor\}.
\]
The correction column of \(\widehat G_{14}\) is
\(\gamma=\AppThirteenFourteenGamma\), whose projective class is this unique
element; rebuilding recovers \(\widehat G_{14}\).
\end{prop}

\begin{proof}
Direct calculation gives \(G_{14}G_{14}^{\mathsf T}=0\) and
\(\rank G_{14}=7\). Every one of its
\(\binom{14}{7}=3432\) seven-column submatrices has rank~7. Hence every
seven columns are an information set, so the generated code is MDS and has
parameters \([14,7,8]\).  In particular, its minimum-weight coefficient is
the MDS value
\[
A_8=(13-1)\binom{14}{8}=36036.
\]

For the asserted row equivalence, if \(G_{14}^{\sigma}\) denotes the stated
column reordering, then
\[
T=\begin{pmatrix}
12&8&12&5&8&10&0\\
10&11&11&4&3&6&0\\
4&7&2&3&10&9&0\\
9&7&0&9&0&0&0\\
0&0&9&0&7&9&0\\
4&7&0&9&0&5&0\\
0&0&0&0&0&10&7
\end{pmatrix},
\qquad \det T=5,
\]
and exact multiplication gives
\(TG_{14}^{\sigma}=\widehat G_{14}\), with \(\widehat G_{14}\) as in
\eqref{eq:gf13-14-structured-generator}. Applying
\(\partial(s,t)=t-5s\) blockwise gives
\[
\partial\widehat G_{14}=\AppThirteenFourteenReadout.
\]
Thus \(\rank(\partial\widehat G_{14})=6\), while its last row lies in
\(U_5\) and ends in \((1,5)\); hence \(r=1\) and \(D=(1)\).
The universal Gram relation holds for \(P,H,Q\), so
Theorem~\ref{thm:universal-rank-boxed} gives self-duality both before and
after the simultaneous row-and-block deletion.

The parent has \(A_j=0\) for \(1\le j<6\), \(A_6=960\), and
\(A_7=3744\). After projectivizing its weight-six and weight-seven
coefficient vectors, there are \(\AppThirteenFourteenLowLines\) classes.
Exact enumeration of \(\mathbf P^5(\F_{13})\) shows that precisely
\(\AppThirteenFourteenSurvivor\) is nonorthogonal to every one of them.
Finally,
\[
[\AppThirteenFourteenGamma]=\AppThirteenFourteenSurvivor,
\]
and the lower six rows of \eqref{eq:gf13-14-structured-generator}, after
deleting their first block, are literally the parent generator. Reattaching
that block and the first row therefore recovers \(\widehat G_{14}\).
\end{proof}

This calculation supplies a search principle rather than a claim that the
known MDS code itself is new.  Starting from a fixed parent, low-weight
coefficient lines eliminate projective corrections before full minimum
distance is computed.  Here the filter is sharp: exactly one projective class
survives and it reconstructs the optimal child.  For \(m\ge0\), one may
therefore study
\[
\delta_m(G_{14})=\max\{d(C_m): C_{i+1}\text{ is obtained from }C_i
\text{ by a non-direct-sum step of
Lemma~\ref{lem:rank-boxed-building-up}}\}.
\]
The question is whether this fixed-seed search reaches the best possible
minimum distance at length \(14+2m\).  The comparison with known parameters
and the successive constructions below separate the proved examples from this forward search
question.

\FloatBarrier


\subsection{Comparison with known parameters}

For earlier results on self-dual codes over these fields, see
\cite{LeonPlessSloane,BetsumiyaEtAl,KimChoi2021}.  For a displayed generator
\(G\), write
\[
A_j=\#\{a\in\mathbb F_p^k:\operatorname{wt}(aG)=j\}.
\]
Table~\ref{tab:application-distances} compares our exact pair \((d,A_d)\)
with a publicly accessible benchmark at the same length and records every
successive two-coordinate construction used below.  A bold value in our
column matches the smallest \(A_d\) in the publicly accessible comparisons
cited in the table.  At lengths covered by a complete online classification,
or by the MDS weight enumerator, this is a certified minimum.  At unclassified
lengths, bold records only the smallest public coefficient located in this
comparison, not a global minimum claim.  This distinction matters over
\(\GF{13}\): the four published length-12
representatives have different weight enumerators, and exact enumeration of
the generator published by Kim and Choi~\cite{KimChoi2021} gives
\(A_8=1752\) at length~18, so our value (1896), although below the earlier
value (2484), is not bold.  The complete classifications used for the
binary and short quinary rows are available in the online databases of
\href{https://www.math.is.tohoku.ac.jp/~mharada/research/codes/sd2.htm}
{Harada and Munemasa} and
\href{https://www.math.is.tohoku.ac.jp/~mharada/research/codes/sd5.htm}
{Harada and Munemasa}, respectively.  For MDS rows,
\(A_d=(p-1)\binom nd\) is determined by the parameters.

\begin{table}[htbp]
\centering
\scriptsize
\setlength{\tabcolsep}{3pt}
\begin{tabular}{@{}lc>{\centering\arraybackslash}p{77pt}
  >{\centering\arraybackslash}p{87pt}>{\raggedright\arraybackslash}p{121pt}@{}}
\toprule
Code & \(p\) & Our \([n,k,d];A_d\) & Public benchmark \((d_{\rm best};A_d)\) & Reference and construction\\
\midrule
\GolayCatalogueRows
\specialrule{1.1pt}{2pt}{2pt}
\ApplicationCatalogueRowsFive
\GFFiveRepeatedCatalogueRows
\specialrule{1.1pt}{2pt}{2pt}
\ApplicationCatalogueRowsThirteen
\GFThirteenRepeatedCatalogueRows
\bottomrule
\end{tabular}
\caption{Parameters of the constructed codes and public distance benchmarks.}
\label{tab:application-distances}
\end{table}

Each code symbol in Table~\ref{tab:application-distances} links to its exact
generator and verification record.  The numerical verifiers check the
relevant boxed or repeated relation, rank, Gram matrix, and recorded weight
data; every complete distribution is also checked against the MacWilliams
identity.  The three largest repeated generators are written out in full
below.  The finite enumerations are external exact computations rather than
Lean proofs.

\subsection{Largest repeated boxed presentations}

The same two-coordinate border displays the largest code considered over each
of the three fields.  If \(G\) has rows \(g_1,\ldots,g_m\), let
\begin{equation}
\mathcal B_c(G,x)=
\begin{pmatrix}
0&1&x\\
-cy_1&-y_1&g_1\\
\vdots&\vdots&\vdots\\
-cy_m&-y_m&g_m
\end{pmatrix},
\qquad y_i=x\mathbin{\cdot}g_i,\qquad x\mathbin{\cdot}x=-1.
\label{eq:largest-repeated-box}
\end{equation}
For \(q=2,5,13\), take \(c=1,2,5\), respectively.  In characteristic two,
\(-1=1\) and the lower heads become \((y_i,y_i)\), so
\(\mathcal B_1\) is literally the binary repeated box.  With the coordinate
orders recorded in the verification records, the three largest presentations
are
\begin{equation}
\begin{aligned}
\widetilde G_{24}^{(2)}
  &=\mathcal B_1(G_{22}^{(2)},x_{22}^{(2)}),\\
\widetilde G_{24}^{(5)}
  &=\mathcal B_2(G_{22}^{(5)},x_{22}^{(5)}),\\
\widetilde G_{20}^{(13)}
  &=\mathcal B_5(G_{18}^{(13)},x_{18}^{(13)}).
\end{aligned}
\label{eq:three-largest-repeated-boxes}
\end{equation}
Each equality is an exact generator identity after the recorded row
normalization and coordinate permutation, not only an abstract equivalence of
codes.

The extended binary Golay code gives a compact instance in which two exact
levels can be followed in either direction from length~20.  Let
\(\mathcal G_{24}\) be the extended Golay code with its standard generator
\(G_{24}\), and put \(C_{24}=\mathcal G_{24}\).  Repeatedly
apply the two-coordinate reduction of Theorem~\ref{thm:binary-cz-kim} to the
first coordinate pair of the current tail.  Deterministic row reduction gives
self-dual generators \(G_m\), for \(m=16,18,20,22\), such that
\begin{equation}
 \mathcal B_1(G_m,x_m)
 =\begin{pmatrix}1&0&x_m\\ y_m&y_m&G_m\end{pmatrix}
 \quad\text{generates }C_{m+2},
 \qquad m\in\{16,18,20,22\},
 \label{eq:golay-lineage-kim}
\end{equation}
where \(C_m=\operatorname{rowspan}(G_m)\) and
\[
\begin{aligned}
x_{16}&=\GolaySixteenCorrection,&
x_{18}&=\GolayEighteenCorrection,\\
x_{20}&=\GolayTwentyCorrection,&
x_{22}&=\GolayTwentyTwoCorrection,\\
(y_m)_i&=x_m\mathbin{\cdot}(G_m)_i.
\end{aligned}
\]
Here \(x_{22}^{(2)}=x_{22}\) in
\eqref{eq:three-largest-repeated-boxes}.
The parent in each of Figures~\ref{fig:binary-largest-repeated}--
\ref{fig:gf13-largest-repeated} is itself put in the universal form
\(G(c;b,\ell,D)\), rather than left as an arbitrary generator.  Within the
parent, the teal rows and paired columns display the forced diagonal blocks
\((a_i,ca_i+1)\) and the off-diagonal blocks \(b_{ij}(1,c)\).  The violet
terminal rows and columns display \(\ell_i\) and the invertible core
\(D_{st}(1,c)\); the opposite block is forced to be
\(-(D\partial_c\ell_i)_s(1,c)\).  Blue is the newly adjoined row and orange is the newly
adjoined coordinate pair.  The outer thick rules mark the building-up border,
whereas the inner thick rules separate the pivot part from the terminal part of
the universal parent.  A light gray background marks every diagonal
two-coordinate block, including the new block and the terminal core.
For the binary endpoint the displayed parent has
\(\BinaryLargestParentParameters\).  Adding the all-ones master row when
necessary and then permuting the two coordinates in the corresponding pivot
pair puts every diagonal block in the form (01) and every terminal pivot
block in the form (10).  The last row is (11) in every block.  Hence the
parent is displayed literally in the form of
Theorem~\ref{thm:binary-universal-rank-one}, not merely as a general
rank-one universal box.
\begin{figure}[htbp]
\centering
\resizebox{\textwidth}{!}{$
\widetilde G_{24}^{(2)}=\BinaryLargestRepeatedMatrix .
$}
\caption{Binary repeated realization \(C_{22}^{(2)}\longleftrightarrow
\mathcal G_{24}\) in the form of Theorem~\ref{thm:binary-universal-rank-one}.}
\label{fig:binary-largest-repeated}
\end{figure}
\FloatBarrier
Thus length~20 is the center of the requested two-level picture:
\[
C_{16}\longleftarrow C_{18}\longleftarrow C_{20}
\longrightarrow C_{22}\longrightarrow C_{24}=\mathcal G_{24}.
\]
Following the arrows to the left means two-coordinate reduction; the right
arrows are Kim building-up with \(x_{20}\) and \(x_{22}\).  Every left edge is
also inverted exactly by the corresponding Kim step with \(x_{16}\) or
\(x_{18}\), so all four adjacent pairs satisfy the same reconstruction
identity.
The parameters, minimum-word counts, and exact upward relations at every level
are recorded in Table~\ref{tab:application-distances}.  In particular, the
distance rises at both upward steps from the length-20 center.  All five
distances attain the
classical upper bounds for binary self-dual codes~\cite{HuffmanPless}.

Over \(\GF{5}\), let \(C_{24}^{(5)}\) be the self-dual
\([24,12,9]\) code of Georgiou and Koukouvinos~\cite{GeorgiouKoukouvinos}.
Its exact reduction at the ordered coordinate pair \((1,14)\) gives the
self-dual \([22,11,8]\) parent \(C_{22}^{(5)}\).  For the displayed parent
generator \(G_{22}^{(5)}\),
\[
x_{22}^{(5)}=\GFFiveTwentyTwoCorrection
\]
has norm \(-1\), and
\[
\operatorname{rowspan}\mathcal B_2
  (G_{22}^{(5)},x_{22}^{(5)})=C_{24}^{(5)}.
\]
In the same paired coordinate order, the complete normalized generator is
the parent--child figure; here the universal parent has
\(\GFFiveLargestParentParameters\).  Thus it is literally the rank-one split
box of Corollary~\ref{thm:split-boxed-form}: every diagonal block is \(01\),
the off-diagonal blocks are \(b_{ij}(1,2)\), and the final block is \((1,2)\).
\begin{figure}[htbp]
\centering
\resizebox{\textwidth}{!}{$
\widetilde G_{24}^{(5)}=\GFFiveLargestRepeatedMatrix .
$}
\caption{Repeated realization over \(\GF{5}\), with the parent in the form of
Corollary~\ref{thm:split-boxed-form}.}
\label{fig:gf5-largest-repeated}
\end{figure}
\FloatBarrier
The checker verifies the literal matrix in
\eqref{eq:three-largest-repeated-boxes}, its parent row space, and the complete
weight distribution of \(C_{22}^{(5)}\); the exhaustive Magma audit finds
132 distance-nine-loss-free ordered pairs among all \(24\cdot23=552\), and
every resulting parent has \(A_8=660\).

Over \(\GF{13}\), let \(C_{20}^{(13)}\) be the pure double-circulant
self-dual \([20,10,10]\) code whose circulant first row is
\[
(3,7,5,5,2,8,5,6,3,0),
\]
as in~\cite{BetsumiyaEtAl}.  Exact two-coordinate reduction at the ordered
pair \((11,7)\) gives a self-dual \([18,9,8]\) code
\(C_{18}^{(13)}\).  The coordinate indices are one-based, and
\GFThirteenRepeatedCertificate\ links to the matrices and their verification
record.  Hence a single reduction--reconstruction pair attains the published
best-known distances at both lengths.

For the displayed universal parent generator \(G_{18}^{(13)}\), the correction
vector is
\[
x_{18}^{(13)}=\GFThirteenEighteenCorrection.
\]
It satisfies \(x_{18}^{(13)}\mathbin{\cdot}x_{18}^{(13)}=-1\) and the exact
row-space identity
\[
\operatorname{rowspan}\mathcal B_5
  (G_{18}^{(13)},x_{18}^{(13)})=C_{20}^{(13)}.
\]
The complete normalized generator at the largest level is
formed from the universal parent with
\(\GFThirteenLargestParentParameters\).  It likewise has the literal
Corollary~\ref{thm:split-boxed-form} pattern, with diagonal blocks \(01\),
off-diagonal blocks \(b_{ij}(1,5)\), and final block \((1,5)\):
\begin{figure}[htbp]
\centering
\resizebox{\textwidth}{!}{$
\widetilde G_{20}^{(13)}=\GFThirteenLargestRepeatedMatrix .
$}
\caption{Repeated realization over \(\GF{13}\), with the parent in the form of
Corollary~\ref{thm:split-boxed-form}.}
\label{fig:gf13-largest-repeated}
\end{figure}
\FloatBarrier
Consequently the two representatives attaining the published best-known
distances form the bidirectional construction
\[
C_{18}^{(13)}\ \longleftrightarrow\ C_{20}^{(13)}.
\]
The checker verifies the literal repeated matrix, equality with the child row
space, and the complete weight distributions at lengths 18 and 20.  The Magma
audit examines all \(20\cdot19=380\) ordered pairs.  Every pair loses a
weight-ten word, so this particular \([20,10,10]\) code cannot yield a
distance-nine parent by a single two-coordinate reduction; the minimum loss is
1896 words, attained by ten ordered pairs.  This is a limitation of the fixed
code and pairing search, not a nonexistence result for a self-dual
\([18,9,9]\) code.

For every upward edge in Table~\ref{tab:application-distances}, the corresponding
checker verifies rank, self-orthogonality, literal reconstruction of the
displayed repeated matrix, equality with the child row space, and the recorded
weight distribution together with its MacWilliams identity.

\section{Conclusion}
\label{sec:conclusion}

We have studied self-dual codes over finite fields $\F_q$ with $q\equiv 1 \pmod 4$ from three complementary viewpoints: the classical building-up construction, the hyperbolic geometry forced by the split condition $-1\in (K^\times)^2$, and a Lean~4 formalization of the resulting algebraic statements.  Section~\ref{sec:applications} exhibits repeated boxed presentations over all three fields considered here: the binary Golay endpoint \([24,12,8]\), a self-dual \([24,12,9]\) code over \(\GF{5}\), and the best-known-distance sequence \([18,9,8]\leftrightarrow[20,10,10]\) over \(\GF{13}\).  Each is displayed through the same two-coordinate border and an explicit parent.  The universal form also gives an explicit rank-one realization and exact two-coordinate parent of a self-dual MDS \([14,7,8]\) code over \(\GF{13}\).  The algebraic results are formalized in Lean~4.

In the binary case, the universal rank-one normalization proves intrinsically that every Euclidean self-dual code has the literal Chinburg--Zhang boxed representative under a scalar-coordinate permutation.  Theorem~\ref{thm:binary-cz-kim} then identifies the Chinburg--Zhang cohomological pairing with the Euclidean form by a specified orthonormal-basis isometry and compares the resulting boxed matrix with Kim's building-up extension.  The top-down deletion is a coordinate reduction through a non-alternating first coordinate plane, not deletion of a hyperbolic pair.
Over \(\F_q\) with \(q\equiv 1\pmod 4\), the decomposition obtained from
the linear map \(\partial\) in Theorem~\ref{thm:universal-rank-boxed} gives the universal
analogue: every Euclidean self-dual code has a rank-\(r\) boxed
representative after a permutation of two-coordinate blocks.  The intrinsic
integer for the chosen ordered coordinate pairing is
\(r=\dim(C\cap U_c)\), and the
terminal matrix \(D\) is nonsingular, with its relation to the other blocks
specified by the Gram equations.  The explicit split boxed
construction is the rank-one specialization with zero pivot diagonal, while
the adapted reverse theorem records the relations forced by that normalization.
Simultaneous pivot-row and block-column deletion preserves the universal
form, giving a chain of self-dual codes ending at \(U_c^{(r)}\).
Lemma~\ref{lem:rank-boxed-building-up} gives the converse operation with
fixed \(D\), and Lemmas~\ref{lem:repeated-kim-lee}--\ref{lem:kim-lee-repeated}
identify its exact fixed-parent relation to Kim--Lee building-up. The
case of a zero correction column is a direct-sum extension, not a
Kim--Lee step in that fixed coordinate order. These results give a
universal representation and recursive construction, not unique parameters
or a classification of permutation-equivalence classes. They also do not
assert an odd-characteristic arithmetic realization by Hilbert symbols.

The distinction between the bilinear forms is essential. An isometry
from the cohomological pairing to the Euclidean form permits comparison
with coding-theoretic constructions, whereas a change to hyperbolic
coordinates does not imply permutation equivalence. Throughout the
normal-form theorems, coordinate permutations preserve the Euclidean form.


\section*{Availability of the formalization}

The Lean proofs, independent theorem comparisons, and verification
instructions, together with the reproducible application data and scripts,
are provided in the public repository
\url{https://github.com/LeGenAI/intersection-coding-theory-cohomology}.
The version accompanying this manuscript is tagged
\href{https://github.com/LeGenAI/intersection-coding-theory-cohomology/tree/afm-revision-2026-08-27}{\texttt{afm-revision-2026-08-27}}.

\bibliographystyle{ieeetr}
\begingroup
\small
\begin{thebibliography}{99}
\setlength{\itemsep}{3pt}
\setlength{\parskip}{0pt}



\bibitem{BacGab} C. Bachoc and P. Gaborit, ``Designs and self-dual codes with long shadows,'' \textit{J. Combin.
Theory Ser. A,} 2004, 105(1): 15-34.

\bibitem{BaDou}
 E. Bannai, S. T. Dougherty, M. Harada, and M. Oura, Type II codes, ``even unimodular
lattices, and invariant rings,'' \textit{IEEE Trans. Inf. Theory,} 1999, 45(4): 1194-1205.

\bibitem{BetsumiyaEtAl}
K.~Betsumiya, S.~Georgiou, T.~A.~Gulliver, M.~Harada, and C.~Koukouvinos,
``On self-dual codes over some prime fields,''
\textit{Discrete Math.}, vol.~262, nos.~1--3, pp.~37--58, 2003.

\bibitem{CRSS}
A. R. Calderbank, E. M. Rains, P. M. Shor, and N. J. A. Sloane, Quantum error
correction via codes over GF(4), IEEE Trans. Inf. Theory, 1998, 44(4): 1369-1387.


\bibitem{ChinburgZhang}
T.~Chinburg and Y.~Zhang,
``Every binary self-dual code arises from Hilbert symbols,''
\textit{Homology, Homotopy Appl.}, vol.~14, no.~2, pp.~189--196, 2012.


\bibitem{ConSlo}
J. H. Conway, N. J. A. Sloane, Sphere Packings, Lattices and Groups, 3rd edn.
Springer, New York (1998).

\bibitem{GeorgiouKoukouvinos}
S.~Georgiou and C.~Koukouvinos,
``New self-dual codes over \(\GF{5}\),''
in \textit{Cryptography and Coding}, Lecture Notes in Computer Science,
pp.~63--69, Springer, 1999.





\bibitem{HuffmanPless}
W.~C.~Huffman and V.~Pless,
\textit{Fundamentals of Error-Correcting Codes},
Cambridge University Press, 2003.


\bibitem{Kim2001}
J.-L. Kim,
``New extremal self-dual codes of lengths 36, 38, and 58,''
\textit{IEEE Trans.\ Inform.\ Theory}, vol.~47, no.~1, pp.~386--393, 2001.

\bibitem{KimChoi2021}
J.-L. Kim and W.-H.~Choi,
``Self-dual codes, symmetric matrices, and eigenvectors,''
\textit{IEEE Access}, vol.~9, pp.~104294--104303, 2021.

\bibitem{KimLee}
J.-L.~Kim and Y.~Lee, ``Euclidean and Hermitian self-dual MDS codes over large finite fields.'' J. Combin. Theory Ser. A 105(1):79-95, 2004.


\bibitem{KimLee2009}
J.-L. Kim and Y.~Lee,
``Self-dual codes using the building-up construction,''
\textit{IEEE Int.\ Symp.\ Inform.\ Theory}, Seoul, Korea (South), 2009, pp. 2400--2402.


\bibitem{KreckPuppe}
M.~Kreck and V.~Puppe,
``Involutions on 3-manifolds and self-dual, binary codes,''
\textit{Homology, Homotopy Appl.}, vol.~10, no.~2, pp.~139--148, 2008.

\bibitem{LeonPlessSloane}
J.~S.~Leon, V.~Pless, and N.~J.~A.~Sloane,
``Self-dual codes over GF(5),''
\textit{J.\ Combin.\ Theory Ser.~A}, vol.~32, no.~2, pp.~178--194, 1982.

\bibitem{Mon}
H. Montani,
``Lagrangian and orthogonal splittings, quasitriangular Lie bialgebras and almost complex product structures,''
\textit{J. Math. Phys.}, vol.~64, 053505, 2023.

\bibitem{Ball2012}
S.~Ball, ``On sets of vectors of a finite vector space in which every subset
of basis size is a basis,'' \textit{J. Eur. Math. Soc.}, vol.~14, no.~3,
pp.~733--748, 2012. doi:10.4171/JEMS/316.

\bibitem{Pless1965}
V.~Pless,
``The number of isotropic subspaces in a finite geometry,''
\textit{Atti Accad.\ Naz.\ Lincei Rend.}, vol.~39, pp.~418--421, 1965.


\bibitem{self-dual}
 E. Rains and N. J. A Sloane, Self-dual codes, in: V. S. Pless, W. C. Huffman (Eds.),
Handbook of Coding Theory, Elsevier, Amsterdam. The Netherlands. (1998)

\bibitem{Ser}
 J.-P. Serre, A course in arithmetic, Grad. Texts in Math. 7, Springer-Verlag,
New York, 1973.

\bibitem{Sha}
C.~E.~Shannon, ``A mathematical theory of communication,''
\textit{Bell System Technical Journal}, vol.~27, no.~3, pp.~379--423,
and no.~4, pp.~623--656, 1948.


\bibitem{ShiChoShaSol}
 M.~Shi, Y.~J.~Choie, A.~Sharma, and P.~Sol\'{e}, Codes and Modular Forms: A Dictionary. World Scientific, 2020.





\end{thebibliography}
\endgroup

\end{document}
